% !TeX root = ../kazymyr-thesis.tex

\chapter{Experiments}\label{chapter:experiments}

The hybrid mode is evaluated with five parameter sweeps, 640 simulation
runs in total. The first three change the system itself and the last two
change the workload it has to handle:

\begin{itemize}
  \item \textbf{X1 --- reconciler processors}
        (Section~\ref{sec:x1-reconciler-sweep}). Varies how much capacity is
        available for reconciliation. This shows in which range of settings the
        adaptive oracle switches modes at all, and what running a mixture of the
        two modes costs compared with running either mode alone.

  \item \textbf{X2 --- critical threshold}
        (Section~\ref{sec:x2-critical-sweep}). Varies the threshold $c$ at which
        the oracle turns ACID on. This tests whether $c$ behaves as a
        consistency budget, that is, whether the measured inconsistency stays
        below the value that was configured, and what throughput is given up for
        a tighter budget.

  \item \textbf{X3 --- adaptive versus random oracle}
        (Section~\ref{sec:x3-random-baseline}). Replaces the measured dirty
        fraction by a fixed probability: each transaction is assigned a mode by
        a coin flip instead of by the state of the system. Both oracles can
        reach the same ACID share, so the comparison isolates what the feedback
        signal contributes from what simply running some ACID contributes.

  \item \textbf{X4 --- semantic overlap}
        (Section~\ref{sec:x4-overlap-sweep}). Varies how much of the system one
        write reaches, by changing the fraction of resources tied to resources
        in other models, and therefore how much work one write creates
        elsewhere.

  \item \textbf{X5 --- resources per transaction}
        (Section~\ref{sec:x5-resources-sweep}). Varies how many resources a
        single transaction touches, and therefore how deep one unit of
        submitted work reaches into the schema.
\end{itemize}

The five sweeps answer RQ2 and RQ3. RQ1 is not among them, and it is worth
saying why rather than leaving the reader to notice the omission. RQ1 asks
how a hybrid mode can be integrated safely into a multi-model transactional
system, which is a design question: it is settled by constructing something
that works and arguing that it is correct, not by varying a parameter and
measuring the result. Its answer is therefore Chapter~\ref{chapter:approach}
--- the mixing hazard, the locked-XOR-dirty rule that resolves it, and the
argument that the rule adds no deadlock --- and no sweep can confirm or
refute it in the way a sweep settles a question about throughput or
inconsistency.

What experiments can do is show that the design behaves as claimed once it
is running. Section~\ref{sec:x0-arbiter} does this before the sweeps begin:
it reports how often the mixed state was actually reached during the 640
runs and resolved by the rule, which establishes that the guarantee is not
vacuous. It is a validation of the construction rather than a sweep, which
is why it carries no parameter of its own. Every later section then depends
on it, because a performance number measured under a rule that was not well
defined would mean nothing.

The remaining questions map onto the sweeps directly. RQ2, which compares
the hybrid mode with pure ACID and pure BASE and asks whether the
inconsistency it permits can be configured, is answered by X1 and X2, with
X4 and X5 establishing the range of workloads over which that configurability
holds. RQ3, which asks whether choosing the mode from the system's own
metrics beats a random assignment, is answered by X3.

\section{Experimental settings}\label{sec:exp-setup}

\subsection{Simulator baseline configuration}\label{sec:exp-baseline}

Every experiment varies one or two parameters and holds everything else at
the baseline of Table~\ref{tab:exp-baseline}.

\begin{table}[htbp]
  \centering
  \caption{Baseline configuration.}
  \label{tab:exp-baseline}
  \small
  \begin{tabular*}{\textwidth}{@{\extracolsep{\fill}} l l l}
    \toprule
    Parameter                        & Baseline value       & What it controls                                  \\
    \midrule
    \multicolumn{3}{l}{\emph{System}} \\
    Resources $N$                    & 5000                 & Size of the data set                              \\
    Models                           & 10                   & Number of co-existing data models                 \\
    Semantic overlap $\sigma$        & 0.30                  & How much information the models share (X4)        \\
    IO processors                    & 30                   & Capacity for IO work                              \\
    CPU processors                   & 30                   & Capacity for CPU work                             \\
    Locks                            & 5000                 & Available locks                                   \\
    Reconcilers $R$                  & swept                & Processors reserved for reconciliation (X1)       \\
    CPR strategy                     & \texttt{ON\_CASCADE} & When cross-model repairs are triggered            \\
    \texttt{maxBookkeeperDepth}      & 5                    & How far a repair may cascade                      \\
    \midrule
    \multicolumn{3}{l}{\emph{Workload}} \\
    Profile                          & \texttt{EQUILIBRIUM} & Keeps the number of in-flight transactions fixed  \\
    In-flight transactions           & 30                   & The constant that the profile holds               \\
    Resources per transaction $L$    & 50                   & Most resources one transaction touches (X5)       \\
    Models per transaction           & up to 4              & Models one transaction may span                   \\
    \midrule
    \multicolumn{3}{l}{\emph{Oracle}} \\
    Critical threshold $c$           & 0.10                 & Dirty fraction at which ACID is switched on (X2)  \\
    Recovery threshold $r$           & 0.05                 & Dirty fraction at which BASE is resumed (Sec.~\ref{sec:effective-recovery}) \\
    \midrule
    \multicolumn{3}{l}{\emph{Run}} \\
    Duration                         & 50\,000\,s           & Simulated time per run                            \\
    Seed pairs                       & 5                    & Replicates per configuration                      \\
    \bottomrule
  \end{tabular*}
\end{table}

The oracle compares the two thresholds against the dirty fraction as it is
at the moment of the decision; there is no filter and no further tuning
parameter (Section~\ref{sec:alpha-ablation}). The two are not realised
equally well: $c$ is met closely, while $r$ turns out to be an upper bound
on where an ACID excursion ends rather than the point at which it ends.
That asymmetry applies to every experiment below, so it is reported once,
at the end of the chapter (Section~\ref{sec:effective-recovery}).

The simulator generates its own workload, so no external data set is
involved. The \texttt{EQUILIBRIUM} profile keeps the number of in-flight
transactions constant: whenever a transaction completes, a new one is
admitted. The load is therefore stationary by construction, which is what
makes throughput comparable across configurations, and also what prevents
this chapter from testing adaptation to a load that changes over time
(Section~\ref{sec:exp-threats}).\label{sec:exp-workload}

One value deviates from the reference configuration of the base simulator:
a run lasts 50\,000\,s of simulated time rather than 10\,000\,s. The
reason is that switching modes is a slow event. The oracle only looks at the
dirty fraction when a new transaction is admitted, and by the time the
system enters ACID mode a large share of the resources is already dirty, so
the ACID transactions that follow have to wait for resources that are still
being repaired. Getting back out of ACID mode is then limited by how much
capacity is available for reconciliation, which at the low end of the sweep
is a single processor. A run therefore stays in one mode for long stretches,
and a 10\,000\,s run would contain only a handful of switches --- too few
to read a stable 95th percentile from. At 50\,000\,s every run in the range
where the oracle is active contains enough of them for the tail to be
meaningful.

\subsection{Implementation and execution}\label{sec:exp-implementation}

Each experiment is driven by its own runner script. The script takes the
parameters that the experiment varies, runs every combination of them under
five seed pairs, and writes one row per run into a CSV file. The exported result files
and the analysis scripts behind the tables in this chapter are part of the
supplemental material.

Two seeds define a system instance. The \emph{semantic-overlap seed} fixes
which resources of which models represent the same information, and the
\emph{transaction-sampling seed} fixes the arrival sequence: how many
resources each transaction touches, which ones, and in which models. The
same five seed pairs are used for every mode in every experiment, so ACID,
BASE and both HYBRID configurations see exactly the same transactions on
exactly the same overlap structure; the pairs themselves are listed with the
supplemental material. All tables report the mean over the five
replicates.\label{sec:exp-seeds} Because the simulator's clock is virtual
and all seeds are fixed, a run reproduces exactly, independently of the
host machine.

\subsection{Metrics}\label{sec:exp-metrics}

Three numbers carry the argument, and each of them is reported in almost
every table of this chapter.

\begin{description}
  \item[Throughput] (\texttt{throughput\_tps}) --- root transactions
        completed per second of simulated time. This is the performance side of
        the trade-off: keeping the data consistent is only useful if the system
        still does work.

  \item[$p_{50}$ dirty fraction] (\texttt{dirty\_p50\_pct}) --- the median
        share of resources that are inconsistent. This describes the typical
        state of the system.

  \item[$p_{95}$ dirty fraction] (\texttt{dirty\_p95\_pct}) --- the 95th
        percentile of the same quantity, that is, the value that inconsistency
        stays below for 95\% of the simulated time. This describes the bad case,
        and it is the number a threshold is supposed to bound: a mixture of
        modes can hit any average, so only the tail shows whether the budget
        actually holds.
\end{description}

Both percentiles are taken over samples spaced evenly in simulated time, so
they are statements about time and not about transactions. Sampling happens
after the work of each tick, so a reading is always a post-repair value and
short peaks inside a tick are not seen. Pure ACID runs have no
dirty-fraction data at all, since inconsistency cannot occur there by
construction; those cells read ``n/a by construction''.

Two further numbers say how much ACID a run actually performed, and they
are not interchangeable:

\begin{description}
  \item[$\phi_{\mathrm{time}}$] (\texttt{acid\_fraction\_time}) --- how much
        of the run the system spent in ACID mode. The oracle timestamps every
        decision and attributes the interval since the previous one to the mode
        that was in force across it, so the total ACID time is known exactly;
        divided by the run duration it gives $\phi_{\mathrm{time}}$. The
        accounting is the same code for both oracles, so the column is
        comparable between them.

  \item[$\phi_{\mathrm{comp}}$] (\texttt{acid\_fraction\_completed}) --- how
        much of the finished work ran under ACID: the number of root
        transactions completed in ACID mode, divided by all completed root
        transactions.
\end{description}

The two can differ widely, because one ACID transaction occupies far more
time than one BASE transaction.


% TODO decide where this belongs
\todo{decide where this belongs} Bookkeeper
transactions are not counted in $\phi_{\mathrm{comp}}$. They inherit the
mode of the root transaction that created them, so counting them would
weight each mode by how many repairs it happened to trigger rather than by
the work that was actually submitted.

% TODO the two different samplers should have proper explanation
\todo{two different samplers -> own chapter}
Both of these numbers, and the dirty fractions above, come from reading the
state of the system while it runs --- but not every reading is taken by the
same reader, and that distinction is what
Section~\ref{sec:effective-recovery} turns on. The \emph{signal} sampler
is the oracle's sensor: it is reached only when a root transaction arrives
without a mode tag, so its estimate advances once per admission, on no
timer, at whatever rate the workload happens to produce. The \emph{metrics}
sampler is the measuring instrument: once per simulated second, regardless
of mode. The sensor is therefore load-gated --- its sampling rate is
roughly the throughput that its own decisions suppress. Clocking the
instrument the same way would make the two indistinguishable, since an
apparent shift in the operating point could equally be the instrument
slowing down together with the system it measures.

Two further columns are used wherever the load on reconciliation is at
issue. \texttt{dirty\_mark\_events\_total} counts transitions of a
resource from clean to dirty, so over the run duration it is the rate at
which repair work is created; \texttt{dirty\_clear\_events\_total}
counts repairs completed, so over the duration it is the rate at which
that work is served, against a ceiling of $0.5R$ resources per second
(Section~\ref{sec:recovery-mechanism}). A write to a resource that is
already dirty is counted in \texttt{dirty\_mark\_calls\_total} but
creates no event and no additional repair, which is why the two mark
columns differ.

Transaction latency (\texttt{txn\_latency\_mean\_ms}) is the time from a
transaction's creation to its commit, recorded for completed root
transactions only; bookkeepers are excluded, as they are everywhere else.
It is the measure of how long work sits in the system, and X5
(Section~\ref{sec:x5-resources-sweep}) uses it because there a transaction
is not a fixed amount of work: the simulator charges processing time per
resource touched, 200\,ms of IO and 100\,ms of CPU each, so a larger
transaction occupies the system proportionally longer.


\section{X0: Arbiter validation}\label{sec:x0-arbiter}

RQ1 asks how a hybrid mode can be integrated \emph{safely}. The answer is
given by construction in Section~\ref{sec:approach-arbiter}: the
locked-XOR-dirty invariant holds because neither write path can produce a
violating state. An ACID transaction cannot acquire a lock on a dirty
resource, and a BASE transaction cannot mark a locked resource dirty; each
defers instead. What construction alone does not show is that the situation
ever arises. A rule that is never reached is trivially never violated, and
the guarantee would then be empty. X0 reports the instrumentation that
distinguishes the two.

Both deferral paths are counted per run. The counter
\texttt{acid\_txn\_blocked\_on\_dirty\_count} records how often an ACID
transaction was made to wait for reconciliation to clear a resource it
needed, and \texttt{base\_txn\_yielded\_on\_locked\_count} how often a BASE
transaction abandoned an update because the resource was locked. Every
increment is one occurrence of the mixing hazard of
Section~\ref{sec:approach-mixing-hazard}, resolved by the rule. X0 is not a
sweep of its own: it aggregates these two columns over all 640 runs of the
five sweeps reported in this chapter.

\begin{table}[htbp]
  \centering
  \caption{X0: arbiter activity across all runs in this chapter, by
    execution mode.}
  \label{tab:x0-arbiter}
  \small
  \begin{tabular}{l r r r}
    \toprule
    Mode             & Runs & ACID waits on dirty & BASE yields on locked \\
    \midrule
    ACID             & 35   & 0                   & 0                     \\
    BASE             & 35   & 0                   & 0                     \\
    HYBRID\_ADAPTIVE & 435  & 550\,793            & 4\,684\,097           \\
    HYBRID\_RANDOM   & 135  & 121\,021            & 1\,078\,556           \\
    \midrule
    Total            & 640  & 671\,814            & 5\,762\,653           \\
    \bottomrule
  \end{tabular}
  \datasource{Columns \texttt{acid\_txn\_blocked\_on\_dirty\_count} and
    \texttt{base\_txn\_yielded\_on\_locked\_count}, summed over all CSVs of
    the five sweeps (X1--X5), five seeds each.}
\end{table}

\begin{description}
  \item[A1] \textbf{The counters are zero in every pure run.} All 35 pure
        ACID and all 35 pure BASE runs record zero on both counters, without
        exception. The arbiter is inert exactly where only one set of
        semantics is in use, which is what confirms that the counters measure
        cross-mode interaction and not ordinary lock or CPR activity.

  \item[A2] \textbf{The hazard is the normal case, not a corner case.}
        Across the hybrid runs the rule fired 6.4 million times, 5.8 million
        of them on the BASE side. The state that neither pure mode defines is
        therefore reached constantly under mixed execution, and the invariant
        holding is a claim with content rather than a claim about a situation
        that never occurred.

  \item[A3] \textbf{Arbiter activity coincides exactly with oracle
          engagement.} Of the 570 hybrid runs, 102 record zero on both
        counters, and in every one of those runs \texttt{mode\_switch\_count}
        is also zero; the correspondence is exact in both directions, with no
        exceptions among the 570. These runs are the configurations
        independently identified as degenerate elsewhere in this chapter:
        $R = 0$ and $R = 25$ in X1 (Section~\ref{sec:x1-engagement}), $c$
        above the BASE steady-state dirtiness in X2
        (Section~\ref{sec:x2-degenerate}), and overlap below the engagement
        threshold in X4 (Section~\ref{sec:x4-engagement}). Degeneracy is thus
        visible in a quantity derived from the arbiter alone, independently of
        the switch count used to establish it elsewhere.
\end{description}

\paragraph{What this does not establish.} X0 is evidence for the invariant,
not for the deadlock-freedom argument of
Section~\ref{sec:approach-deadlock}, which remains analytical. Every run in
this chapter ends by reaching the fixed simulated horizon of 50\,000
seconds rather than by quiescing, so termination is imposed by the harness
and a stalled run would end the same way; the run data cannot separate the
two. The weaker statement it does support is that no run stopped making
progress: the least productive hybrid run in the whole set still committed
336 transactions, and every configuration retains nonzero throughput in
Tables~\ref{tab:x1-summary} and~\ref{tab:x2-tracking}. A direct test would
need a wait-for-graph check at runtime, which is left to future work.

\section{X1: Reconciler-processor sweep}\label{sec:x1-reconciler-sweep}

X1 answers the first half of RQ2: how the hybrid mode compares with pure
ACID and pure BASE in throughput and inconsistency. It is the only sweep in
which all three run under identical conditions. Before that comparison means
anything, two things must be settled: whether the oracle engages at all, and
what the mixture costs when it does. The other half of RQ2, configurability,
is left to X2; X1 holds $c = 0.10$ throughout.

The sweep covers $R \in \{0,1,5,10,15,20,25\}$ across four configurations
--- pure ACID, pure BASE, \texttt{HYBRID\_ADAPTIVE} and
\texttt{HYBRID\_RANDOM} --- and five seeds, 140 runs.

One point about $R$ has to be kept in mind when reading the table.
Reconcilers are not extra hardware: the simulator has a fixed pool of
processors, and every reconciler is taken out of it. At $R = 1$, one
processor repairs and 29 are left for foreground transactions; at $R = 15$
the pool is split in half, 15 repairing and 15 running transactions. Raising
$R$ therefore buys faster repair by giving up exactly the capacity that
produces the work needing repair, which is why throughput falls off towards
the right-hand end of the sweep. $R = 0$ is not a smaller version of the
same setting but a different arrangement altogether, and is discussed on its
own below.

Table~\ref{tab:x1-summary} reports one row per configuration and reconciler
count. Two of its columns were introduced in Section~\ref{sec:exp-metrics}:
$\phi_{\mathrm{time}}$ is the share of the run spent in ACID mode and
$\phi_{\mathrm{comp}}$ the share of completed transactions that ran under
ACID. \emph{Switches} counts how often the oracle changed the mode during
the run, so a row with zero switches is degenerate --- the oracle never
engaged and the run is a pure BASE run under a different label.

\begin{table}[!htb]
  \centering
  \caption{X1: results by reconciler count $R$ and configuration.}
  \label{tab:x1-summary}
  \small
  \begin{tabular}{r l r r r r r r}
    \toprule
    $R$ & Mode             & Thr.\ (tps) & $\phi_{\mathrm{time}}$ & $\phi_{\mathrm{comp}}$ & Switches & $p_{50}$ (\%)                           & $p_{95}$ (\%) \\
    \midrule
    --  & ACID (all $R$)   & 0.0418      & 1.0000                 & 1.0000                 & 0.0      & \multicolumn{2}{c}{n/a by construction}                 \\
    \midrule
    0   & BASE             & 1.9536      & 0.0000                 & 0.0000                 & 0.0      & 3.17                                    & 3.87          \\
    1   & BASE             & 3.8685      & 0.0000                 & 0.0000                 & 0.0      & 29.50                                   & 29.64         \\
    5   & BASE             & 3.3344      & 0.0000                 & 0.0000                 & 0.0      & 27.20                                   & 27.53         \\
    10  & BASE             & 2.6681      & 0.0000                 & 0.0000                 & 0.0      & 23.38                                   & 23.93         \\
    15  & BASE             & 2.0013      & 0.0000                 & 0.0000                 & 0.0      & 17.74                                   & 18.47         \\
    20  & BASE             & 1.3343      & 0.0000                 & 0.0000                 & 0.0      & 8.10                                    & 8.93          \\
    25  & BASE             & 0.6665      & 0.0000                 & 0.0000                 & 0.0      & 0.16                                    & 0.50          \\
    \midrule
    0   & HYBRID\_ADAPTIVE & 1.9536      & 0.0000                 & 0.0000                 & 0.0      & 3.17                                    & 3.87          \\
    1   & HYBRID\_ADAPTIVE & 0.0726      & 0.7776                 & 0.0918                 & 79.8     & 6.80                                    & 10.12         \\
    5   & HYBRID\_ADAPTIVE & 0.3453      & 0.5068                 & 0.0966                 & 234.0    & 5.67                                    & 10.22         \\
    10  & HYBRID\_ADAPTIVE & 0.5883      & 0.2886                 & 0.0694                 & 224.0    & 3.22                                    & 9.98          \\
    15  & HYBRID\_ADAPTIVE & 0.5941      & 0.1667                 & 0.0716                 & 162.4    & 0.43                                    & 9.36          \\
    20  & HYBRID\_ADAPTIVE & 1.3089      & 0.0025                 & 0.0009                 & 3.6      & 8.08                                    & 8.92          \\
    25  & HYBRID\_ADAPTIVE & 0.6665      & 0.0000                 & 0.0000                 & 0.0      & 0.16                                    & 0.50          \\
    \midrule
    0   & HYBRID\_RANDOM   & 0.0849      & 0.7887                 & 0.4968                 & 2137.2   & 0.00                                    & 0.06          \\
    1   & HYBRID\_RANDOM   & 0.0831      & 0.7890                 & 0.4974                 & 2090.8   & 1.02                                    & 2.98          \\
    5   & HYBRID\_RANDOM   & 0.0849      & 0.7873                 & 0.4971                 & 2135.8   & 0.00                                    & 0.18          \\
    10  & HYBRID\_RANDOM   & 0.0850      & 0.7761                 & 0.4969                 & 2137.6   & 0.00                                    & 0.09          \\
    15  & HYBRID\_RANDOM   & 0.0876      & 0.6601                 & 0.4972                 & 2202.0   & 0.00                                    & 0.07          \\
    20  & HYBRID\_RANDOM   & 0.0938      & 0.5751                 & 0.4978                 & 2350.8   & 0.00                                    & 0.07          \\
    25  & HYBRID\_RANDOM   & 0.0987      & 0.5197                 & 0.5003                 & 2471.2   & 0.00                                    & 0.07          \\
    \bottomrule
  \end{tabular}
\end{table}

\paragraph{Why $R = 0$ behaves differently.} The $R = 0$ row looks
contradictory at first: no reconcilers, and yet the
lowest dirty fraction of the whole BASE block together with a throughput in
the middle of the range. Neither follows from repairs being skipped, because
at $R = 0$ they are not skipped. The simulator still creates the same repair
work; it simply has nowhere separate to put it, so the repairs are queued on
the ordinary IO and CPU processors and run there, after the transaction that
caused them has already committed
(Section~\ref{sec:prelim-engine}). Reconciliation at $R = 0$ is not absent,
it is unpartitioned.

The difference the partition makes is not the order of work inside a repair
but how many repairs can be under way at once. Every repair runs its IO
phase and then its CPU phase, in that order, under both arrangements; the
two never overlap within one repair.

\begin{itemize}
  \item \textbf{Dedicated reconcilers ($R \ge 1$).} One worker owns one
        repair until it finishes, advancing it by a single phase per simulated
        step. At most $R$ repairs are in flight, whatever the rest of the pool
        is doing, and each worker finishes one repair every two seconds.

  \item \textbf{No partition ($R = 0$).} Repairs join the ordinary IO and
        CPU queues, so many different repairs are in flight at once across the
        30 processors --- one repair's IO on one processor while another
        repair's CPU runs on another. The number in flight is bounded by the
        pool rather than by $R$.
\end{itemize}

Two consequences follow. Repair at $R = 0$ is faster, not slower: its rate is
limited only by how much of the pool the foreground leaves idle, which in
this workload is more than a handful of dedicated reconcilers would provide,
so the dirty fraction is correspondingly low.

Throughput is nevertheless not the highest, because for $R \ge 1$ it is set
by the processors left over after the carve-out rather than by repair capacity at all --- across the BASE rows it
is proportional to $30 - R$ to within 0.1\,\%, so reconcilers cost
throughput by occupying processors, not by competing for service time. At
$R = 0$ no processors are lost, but foreground work and repair work share
one pool, so each transaction takes roughly twice as long and throughput
lands in the middle of the range. The $R = 0$ point is therefore not the
left endpoint of the $R \ge 1$ curve; it is a different allocation regime
plotted on the same axis, and the two should not be read as one trend.


\subsection{Validation against the reference results}\label{sec:x1-validation}

The BASE line reproduces the result of the base
simulator~\cite{zacouris2026mxs}. Throughput and
inconsistency move from $(1.95\,\text{tps}, 3.17\%)$ at $R = 0$ to
$(3.87, 29.50\%)$ at $R = 1$, then decline monotonically to
$(0.67, 0.16\%)$ at $R = 25$. This is the same shape reported there:
inline CPRs give the lowest inconsistency at capped throughput, a single
dedicated reconciler raises throughput while inconsistency rises sharply,
and many reconcilers drive inconsistency back down at a throughput cost,
with foreground execution starved at the top of the range.

\subsection{Where the oracle engages}\label{sec:x1-engagement}

The adaptive oracle has three regimes rather than one monotone trade-off.
At $R = 0$ it is inert and degenerates exactly to BASE. Between $R = 1$
and $R = 15$ it engages substantially, with $\phi_{\mathrm{time}}$ from
0.17 to 0.78 and 80 to 234 switches; this band contains the throughput
valley. From $R = 20$ upwards it is inert again and throughput free-rides
back to the BASE level. At $R = 0$ and $R = 25$ every
\texttt{HYBRID\_ADAPTIVE} statistic equals the corresponding BASE row in
all five seeds, so those points are excluded from adaptive-versus-BASE
comparisons or marked as degenerate.

\begin{description}
  \item[O1] \textbf{Reconciler count has no effect under pure ACID.}
        Throughput (0.0418\,tps) and the started and aborted transaction counts
        are identical across all seven values of $R$, confirming that the
        reconciliation stage is never invoked under ACID. The ACID row is a flat
        reference line rather than an unswept baseline.

  \item[O2] \textbf{The adaptive mode degenerates to BASE at both ends, and
          the throughput curve between them is a valley}
        (Figure~\ref{fig:x1-o2-valley}). At $R = 0$ dirtiness
        never accumulates enough to cross $c$; at $R = 25$ repair is fast enough
        that it never reaches $c$ either (Figure~\ref{fig:x1-o2-dirty}). In between, throughput collapses from
        1.95\,tps at $R = 0$ to 0.073\,tps at $R = 1$, recovers through
        $R = 5, 10, 15$, and is back at the BASE level by $R = 25$. The two ends
        look alike in the table and have different causes: at low $R$ the oracle
        correctly detects a problem that drains slowly, and pays for it, while
        at high $R$ it simply never fires and the BASE throughput comes for
        free. Adaptivity is confined to an interior band of $R$.
\end{description}

\begin{figure}[H]
  \centering
  \includegraphics[width=\textwidth]{x1_o2_throughput_valley.pdf}
  \caption{O2: the adaptive throughput valley.}
  \label{fig:x1-o2-valley}
\end{figure}

\begin{figure}[H]
  \centering
  \includegraphics[width=\textwidth]{x1_o2_dirty_vs_threshold.pdf}
  \caption{O2: the uncontrolled dirty fraction against $c$, which decides
    whether the oracle fires at all. Dotted lines are $p_{95}$.}
  \label{fig:x1-o2-dirty}
\end{figure}

\begin{description}
  \item[O3] \textbf{ACID share and switch frequency peak at different $R$}
        (Figure~\ref{fig:x1-o3-share-switches}).
        $\phi_{\mathrm{time}}$ falls monotonically from $R = 1$
        ($\phi_{\mathrm{time}} = 0.7776$) onwards, while the switch count peaks
        later, at $R = 5$ (234 against 80 at $R = 1$). $R = 1$ produces few long
        excursions and $R = 5$ many short ones. More switching does not imply
        more time in ACID.
\end{description}

\begin{figure}[H]
  \centering
  \includegraphics[width=\textwidth]{x1_o3_share_vs_switches.pdf}
  \caption{O3: the ACID time share and the switch count peak at different $R$.}
  \label{fig:x1-o3-share-switches}
\end{figure}

\begin{description}
  \item[O4] \textbf{The random oracle is an effective negative control}
        (Figure~\ref{fig:x1-o4-random}).
        With routing independent of the dirty fraction, throughput is nearly flat,
        0.083 to 0.099\,tps across the entire $R$ range, so the valley in O2
        belongs to the feedback loop and not to reconciler provisioning as such.
        The quantity that stays pinned is $\phi_{\mathrm{comp}}$, at 0.4968 to
        0.5003, which is the coin flip behaving as specified, while
        $\phi_{\mathrm{time}}$ falls from 0.7887 at $R = 0$ to 0.5197 at
        $R = 25$: the random oracle controls its routing share exactly and its
        time share not at all.
\end{description}

\begin{figure}[H]
  \centering
  \includegraphics[width=\textwidth]{x1_o4_random_control.pdf}
  \caption{O4: the random oracle as a negative control, flat where the
    adaptive oracle has a valley.}
  \label{fig:x1-o4-random}
\end{figure}

\subsection{Is the hybrid a weighted average of its two
  modes?}\label{sec:x1-mixing-cost}

If the hybrid mode were nothing more than the two pure modes taken in
turn, its throughput would follow from how long it spends in each. A run
that spends a share $\phi_{\mathrm{time}}$ of its time in ACID and the
rest in BASE would deliver

\begin{equation}
  T_{\mathrm{pred}}(R) \;=\; \phi_{\mathrm{time}} \cdot T_{\text{ACID}}
  \;+\; (1 - \phi_{\mathrm{time}}) \cdot T_{\text{BASE}}(R),
  \label{eq:x1-mixture}
\end{equation}

with $T_{\text{ACID}} = 0.0418$\,tps and $T_{\text{BASE}}(R)$ taken from
the BASE rows of Table~\ref{tab:x1-summary}. This is the natural first
guess, and it is worth writing down because it is exactly what switching
between two modes would give if the switching itself cost nothing.

It costs something. Table~\ref{tab:x1-mixing} puts
Equation~\ref{eq:x1-mixture} next to what the runs actually delivered, and
the two disagree wherever the oracle is active. At $R = 1$ the weighted
average predicts 0.893\,tps and the run delivers 0.073, short by a factor
of 12.3. The gap narrows as $R$ rises --- 4.8$\times$ at $R = 5$,
3.2$\times$ at $R = 10$, 2.8$\times$ at $R = 15$ --- but it never closes.
The hybrid stays ahead of pure ACID throughout, so mixing is not simply
the worst of both; it lands far below the blend of the two that its own
time split predicts.

\begin{table}[htbp]
  \centering
  \caption{\texttt{HYBRID\_ADAPTIVE} throughput as measured, against the
    weighted average of Equation~\ref{eq:x1-mixture}.}
  \label{tab:x1-mixing}
  \small
  \begin{tabular}{r r r r r r}
    \toprule
    $R$ & $\phi_{\mathrm{time}}$ & BASE (tps) & Predicted & Observed & Gap          \\
    \midrule
    1   & 0.7776                 & 3.868      & 0.893     & 0.073    & 12.3$\times$ \\
    5   & 0.5068                 & 3.334      & 1.666     & 0.345    & 4.8$\times$  \\
    10  & 0.2886                 & 2.668      & 1.910     & 0.588    & 3.2$\times$  \\
    15  & 0.1667                 & 2.001      & 1.675     & 0.594    & 2.8$\times$  \\
    20  & 0.0025                 & 1.334      & 1.331     & 1.309    & 1.0$\times$  \\
    \bottomrule
  \end{tabular}
  \datasource{Derived from Table~\ref{tab:x1-summary}; pure-ACID throughput
    0.0418\,tps from the same sweep.}
\end{table}

Two checks say the shortfall is real rather than an artefact of the
prediction. The first is that the prediction is accurate where it should
be: at $R = 20$ the oracle barely engages ($\phi_{\mathrm{time}} =
  0.0025$) and Equation~\ref{eq:x1-mixture} lands within 2\% of the
measurement. A formula that missed everywhere would be the wrong yardstick;
this one is right when there is almost nothing to mix, which is what makes
its failure elsewhere worth reporting.

The second is that the shortfall does not depend on which oracle does the
mixing. The random oracle shows it too: at $R = 10$ it spends
$\phi_{\mathrm{time}} = 0.7761$ of the run in ACID, which predicts
0.630\,tps against a measured 0.0850\,tps, a gap of 7.4$\times$. That
oracle ignores the dirty fraction entirely, so the cost belongs to running
the two modes together and not to the rule that chooses between them.

What Equation~\ref{eq:x1-mixture} leaves out is the cost of the switch
itself. An ACID transaction admitted while the resources it needs are
still dirty has to wait for them to be repaired
(Section~\ref{sec:approach-blocking-cost}), and that waiting is invisible
to a calculation built only from how long each mode ran.
Section~\ref{sec:x2-mechanism} measures it.

\section{X2: Consistency-budget sweep}\label{sec:x2-critical-sweep}

The setting this experiment is about the critical dirty-fraction
threshold $c$ in the simulator
(Table~\ref{tab:notation}). It is the share of resources allowed to be
inconsistent before the oracle switches the system to ACID: at every
admission the oracle compares the live dirty fraction $d_k$ against $c$,
and enters ACID mode as soon as $d_k \ge c$
(Equation~\ref{eq:enter-acid}). It is expressed as a fraction of all
resources, so $c = 0.10$ means ``switch to ACID once a tenth of the
resources are waiting to be repaired''. It is the setting through which an
operator states how much inconsistency the system may carry.

X1 establishes that the oracle engages and what happens when it does, but
it holds this threshold fixed at $c = 0.10$, so it says nothing about
whether the value of $c$ makes any difference. That is the claim the
thesis rests on: that an operator can set a level of inconsistency and get
it. X2 tests it by varying $c$ directly and asking whether the measured
inconsistency follows.

The sweep covers $c \in \{0.02, 0.05, 0.10, 0.15, 0.20, 0.30\}$ against
$R \in \{1, 5, 10, 15\}$, five seeds, 120 runs, all
\texttt{HYBRID\_ADAPTIVE}. The recovery threshold is held at $r = 0.5c$
throughout, so the hysteresis band scales with the budget and the sweep
stays one-factor. The values $R \in \{0, 20, 25\}$ are excluded because X1
shows the oracle is inert there.

One property of the baseline sets the top of the useful range for $c$. The
semantic overlap is held at $\sigma = 0.30$, so 30\% of the 5000 resources
carry a cross-model rule and the rest have nothing to reconcile and can
never be marked dirty. The dirty fraction therefore cannot get much above
30\% however the system is run, and the measurements agree: the highest
value recorded anywhere in the 120 runs is 29.88\%. That is why
$c = 0.30$ sits at the top of the sweep. It is not a value an operator
would pick but the point at which the budget reaches the limit of the
system, and it is degenerate at every reconciler count for that reason.
X4 (Section~\ref{sec:x4-overlap-sweep}) varies $\sigma$ itself and shows
what happens to the budget when that limit moves.

The reason for varying the reconciler
count $R$ as well as $c$ is that a setting is only useful if it works the
same way on differently provisioned systems. If the same $c$ gave a
different result on a system with one reconciler than on a system with
fifteen, an operator would have to re-tune it for every deployment, and it
would not be a setting in any practical sense.

\subsection{The critical threshold as a control knob}\label{sec:x2-tracking}

\begin{table}[htbp]
  \centering
  \caption{X2: measured $p_{95}$ dirty fraction against the configured
    threshold.}
  \label{tab:x2-tracking}
  \small
  \begin{tabular}{r r r r r r}
    \toprule
    Configured $c$ (\%) & $R = 1$ & $R = 5$ & $R = 10$ & $R = 15$ & Range (pp) \\
    \midrule
    2                   & 2.52    & 2.32    & 2.10     & 1.84     & 0.68       \\
    5                   & 5.43    & 5.36    & 5.14     & 4.71     & 0.72       \\
    10                  & 10.12   & 10.22   & 9.98     & 9.36     & 0.86       \\
    15                  & 14.76   & 14.79   & 14.56    & 14.28    & 0.51       \\
    20                  & 19.52   & 19.54   & 19.44    & deg.     & 0.10       \\
    30                  & deg.    & deg.    & deg.     & deg.     & --         \\
    \bottomrule
  \end{tabular}
\end{table}

Table~\ref{tab:x2-tracking} reports one measured $p_{95}$ dirty fraction
per cell, in percent and averaged over the five seeds. Each number should
be read against the configured value in the first column of the same
row. ``deg.'' marks a
degenerate cell: there the threshold sits above the inconsistency the
system produces on its own, the oracle never switches to ACID, and the
measured value is simply the BASE level at that reconciler count. The last
column is the range across the four reconciler counts in that row --- the
largest measured value minus the smallest --- counting only the cells where
the oracle engaged.

This is the central result of the thesis. The configured threshold is
varied from 2\% to 30\% and the reconciler count from 1 to 15, and in
every cell where the oracle engages the measured $p_{95}$ stays within
0.74 percentage points of the configured $c$ --- not only on average, but
in each of the 95 individual engaged runs. The largest single deviation
anywhere is $-0.74$\,pp. At $c = 0.10$, the twenty individual runs land
between 9.32\% and 10.28\%. Figure~\ref{fig:x2-tracking} plots the same
data as a distance: how far each cell landed from the value it was set to,
in percentage points. In the terms of Equation~\ref{eq:problem}, $c$ behaves as a
bound on the tail that can be configured: it is set in the same units as
the quantity it controls, and the measured value follows it without having
to be re-tuned per configuration.

\begin{figure}[H]
  \centering
  \includegraphics[width=\textwidth]{x2_tracking.pdf}
  \caption{How far the measured $p_{95}$ dirty fraction lands from the
    configured threshold $c$, one line per reconciler count. Zero is a run
    that hit its setting exactly. The band marks the largest deviation seen
    in any single run, $\pm 0.74$\,pp. Degenerate cells are left out.}
  \label{fig:x2-tracking}
\end{figure}

\paragraph{The threshold is a bound, not a hard ceiling.} The measured
value lands close to $c$ but not exactly on it, and it can land on either
side. Reading the grid, the deviation is $+0.52$\,pp at the top-left
corner ($c = 0.02$, $R = 1$) and $-0.72$\,pp at the bottom-right
($c = 0.15$, $R = 15$), crossing zero in between. The two ends have
different causes.

Going over $c$ is a matter of when the oracle looks. There are two ways
work enters the system, and only one of them is a moment at which the mode
is decided. A root transaction is tagged at admission, and that is when
the oracle reads the dirty fraction. Everything else --- the bookkeeper
transactions that a root transaction spawns --- inherits the tag of its
root and never consults the oracle
(Section~\ref{sec:approach-oracle-interface}). So the dirty fraction is
only checked at root admissions, while work that dirties resources keeps
arriving in between. When root admissions are rare, which is what a low
reconciler count produces, the dirty fraction can climb past $c$ before
the oracle gets its next look. \todo{two admissions --- root admissions
  sample the oracle, inherited bookkeeper work does not; add the reference
  here once the main discussion of this is written}

Staying under $c$ has the opposite cause. Once the system is in ACID mode
repair is quick, so by the time the oracle next looks, the dirty set has
usually already fallen well below the recovery threshold $r$, and the run
spends most of its time comfortably under the budget rather than at it.
Section~\ref{sec:effective-recovery} measures this at every transition.

The right way to state the result is therefore as a magnitude,
$|p_{95} - c| \le 0.74$\,pp, and not as an offset in one direction. The
bound is what an operator can rely on; which side of $c$ the run lands on
depends on the configuration, and both sides are small.

\subsection{Findings}\label{sec:x2-findings}\label{sec:x2-degenerate}\label{sec:x2-throughput}

\begin{description}
  \item[O5] \textbf{The deviation moves from overshoot to undershoot as the
          budget grows} (Figure~\ref{fig:x2-tracking}). At every reconciler
        count the difference between the measured $p_{95}$ and the configured
        $c$ falls steadily as $c$ rises. At $R = 1$, 5 and 10 it starts positive
        at the tight budgets and ends negative at the loose ones, crossing zero
        near $c = 0.10$; at $R = 15$ repair is fast enough that it is already
        negative at the tightest budget. The more reconcilers there are, the
        earlier the crossing happens. The two ends have the causes given in
        Section~\ref{sec:x2-tracking}: with a tight budget and slow repair the
        dirty fraction can climb past $c$ between two root admissions, while
        with a loose budget and fast repair the dirty set is already well below
        the threshold by the time the oracle looks again. Neither side is large
        --- every run stays inside the 0.74\,pp bound --- but the sign is
        predictable rather than random, so an operator setting a small $c$
        should expect to land a little above it, and one setting a large $c$ a
        little below.

  \item[O6] \textbf{Setting the threshold too high switches the oracle off.}
        Five of the 24 cells are degenerate: the system is repaired faster than
        it can get dirty enough to reach $c$, so the oracle never fires and the
        run is pure BASE. This happens exactly when $c$ is above the
        inconsistency BASE produces on its own at that reconciler count, that
        is when $c \ge p_{95}^{\text{BASE}}(R)$, which accounts for all five.
        At $R = 10$ BASE alone runs at 23.93\%, so $c = 0.30$ never fires; at
        $R = 15$ BASE alone runs at 18.47\%, so it stops firing one step
        earlier, at $c = 0.20$. The same rule is used again in X4
        (Section~\ref{sec:x4-overlap-sweep}). Degenerate cells have to be
        spotted before any average is taken over a sweep, which is why every
        HYBRID row in this thesis reports its switch count.

  \item[O7] \textbf{Buying more throughput with the budget works, but it is a
          poor exchange} (Table~\ref{tab:x2-throughput},
        Figure~\ref{fig:x2-throughput}). Allowing more inconsistency does let
        more work through, because the system spends less of its time in the
        slow mode. What it does not do is pay back in proportion: ten times the
        budget returns only one to two and a half times the throughput, and the
        whole range stays far under what the same system reaches without the
        oracle at all --- between a fiftieth and a third of the pure BASE
        throughput at the same reconciler count. At $R = 1$ the exchange all but
        fails: throughput dips before it recovers and ends a fifth above where
        it started, 0.1083\,tps against 0.0907\,tps, for ten times the
        inconsistency. So the threshold is a real control, but an operator
        cannot expect to trade inconsistency for speed at anything like a fair
        rate. Section~\ref{sec:x2-mechanism} explains where the throughput
        goes.

        \begin{table}[H]
          \centering
          \caption{X2: mean throughput and mode-switch count against the
            configured threshold.}
          \label{tab:x2-throughput}
          \small
          \begin{tabular}{r r r r r c r r r r}
            \toprule
                     & \multicolumn{4}{c}{Throughput (tps)} &         & \multicolumn{4}{c}{Mode switches}                                                         \\
            \cmidrule(lr){2-5}\cmidrule(lr){7-10}
            $c$ (\%) & $R{=}1$                              & $R{=}5$ & $R{=}10$                          & $R{=}15$ &  & $R{=}1$ & $R{=}5$ & $R{=}10$ & $R{=}15$ \\
            \midrule
            2        & 0.0907                               & 0.2551  & 0.3951                            & 0.3793   &  & 302     & 509     & 443      & 238      \\
            5        & 0.0718                               & 0.3029  & 0.4860                            & 0.4769   &  & 152     & 347     & 304      & 198      \\
            10       & 0.0726                               & 0.3453  & 0.5883                            & 0.5941   &  & 80      & 234     & 224      & 162      \\
            15       & 0.0840                               & 0.3973  & 0.6895                            & 0.7781   &  & 51      & 162     & 179      & 131      \\
            20       & 0.1083                               & 0.5080  & 0.9460                            & deg.     &  & 36      & 108     & 124      & 0        \\
            30       & deg.                                 & deg.    & deg.                              & deg.     &  & 0       & 0       & 0        & 0        \\
            \bottomrule
          \end{tabular}
          \datasource{Same source as Table~\ref{tab:x2-tracking}; columns
            \texttt{throughput\_tps} and \texttt{mode\_switch\_count}.
            ``deg.'' is a degenerate cell, and a switch count of zero marks
            the same cells.}
        \end{table}

  \item[O8] \textbf{A larger budget makes the oracle switch mode less often}
        (Table~\ref{tab:x2-throughput}). Switch counts fall as $c$ rises at
        every reconciler count. A higher threshold is crossed less often, so
        the run is made of fewer and longer ACID periods rather than many short
        ones.
\end{description}

Both halves of Table~\ref{tab:x2-throughput} move in one direction and neither
moves far. Reading a throughput column downwards, every reconciler count gains
from a larger budget except $R = 1$, which stays flat, and the gain is small
next to the change in the setting that produced it, which is 15 times
from end to end
(Figure~\ref{fig:x2-throughput}). Reading a switch column downwards, the
oracle intervenes steadily less often as the budget grows.

\begin{figure}[H]
  \centering
  \includegraphics[width=\textwidth]{x2_throughput.pdf}
  \caption{O7: throughput against the configured threshold $c$, one line per
    reconciler count. Degenerate cells are left out, which is why the
    $R = 15$ line stops at $c = 0.15$.}
  \label{fig:x2-throughput}
\end{figure}

\subsection{The budget shifts work out of ACID, it does not make ACID
  cheaper}\label{sec:x2-mechanism}

O7 is surprising once Equation~\ref{eq:blocking-probability} is taken into
account. A transaction touches $L = 50$ resources, so its chance of meeting at
least one dirty resource climbs steeply with the dirty fraction: about 64\% at
2\% dirty, 92\% at 5\% and 99.5\% at 10\%. A tight budget should therefore be
the cheap setting --- it keeps the system in the region where an ACID
transaction has a real chance of finding everything it needs already clean,
while a loose budget puts it in the region where it almost never does. The
measurements say otherwise: at every reconciler count above one, the tight
budgets are the slow ones.

The reason is that raising $c$ does not make ACID work cheaper. It makes
less of the work ACID. Each individual ACID transaction gets slower as the
budget grows --- there is more dirt to wait on, exactly as the model predicts
--- but it is a shrinking share of what the system finishes, and the run is
carried instead by BASE transactions that never wait on a repair at all.
Table~\ref{tab:blocking} is arranged to show both movements at once: the left
half is the ACID transactions --- how many there are and what each one costs
--- and the right half is what the run as a whole does.

\begin{table}[H]
  \centering
  \caption{X2: what an ACID transaction costs, and how much of the run is
    ACID.}
  \label{tab:blocking}
  \small
  \begin{tabular}{r r r r c r r r r}
    \toprule
        &      & \multicolumn{2}{c}{ACID transactions} &               & \multicolumn{4}{c}{The run as a whole}                                                                                  \\
    \cmidrule(lr){3-4}\cmidrule(lr){6-9}
    $R$ & $c$  & Adm.\ $A$                             & Mean wait (s) &                                        & $\phi_{\mathrm{time}}$ & $\phi_{\mathrm{comp}}$ & Mean lat.\ (s) & Thr.\ (tps) \\
    \midrule
    1   & 0.02 & 2691                                  & 106.5         &                                        & 0.837                  & 0.472                  & 230.7          & 0.0907      \\
    1   & 0.05 & 1105                                  & 225.9         &                                        & 0.822                  & 0.243                  & 299.1          & 0.0718      \\
    1   & 0.10 & 384                                   & 392.5         &                                        & 0.778                  & 0.092                  & 261.8          & 0.0726      \\
    1   & 0.15 & 199                                   & 538.0         &                                        & 0.734                  & 0.043                  & 208.5          & 0.0840      \\
    1   & 0.20 & 123                                   & 686.5         &                                        & 0.665                  & 0.021                  & 146.4          & 0.1083      \\
    \midrule
    5   & 0.02 & 2788                                  & 18.1          &                                        & 0.440                  & 0.167                  & 78.7           & 0.2551      \\
    5   & 0.05 & 2735                                  & 40.2          &                                        & 0.487                  & 0.147                  & 69.9           & 0.3029      \\
    5   & 0.10 & 1924                                  & 70.2          &                                        & 0.507                  & 0.097                  & 67.2           & 0.3453      \\
    5   & 0.15 & 1186                                  & 95.0          &                                        & 0.481                  & 0.054                  & 60.4           & 0.3973      \\
    5   & 0.20 & 719                                   & 120.8         &                                        & 0.403                  & 0.026                  & 47.8           & 0.5080      \\
    \midrule
    10  & 0.02 & 2805                                  & 7.9           &                                        & 0.271                  & 0.108                  & 50.6           & 0.3951      \\
    10  & 0.05 & 2741                                  & 18.9          &                                        & 0.277                  & 0.088                  & 42.4           & 0.4860      \\
    10  & 0.10 & 2513                                  & 32.5          &                                        & 0.289                  & 0.069                  & 37.4           & 0.5883      \\
    10  & 0.15 & 1979                                  & 43.4          &                                        & 0.286                  & 0.048                  & 33.4           & 0.6895      \\
    10  & 0.20 & 1320                                  & 59.0          &                                        & 0.243                  & 0.024                  & 26.3           & 0.9460      \\
    \midrule
    15  & 0.02 & 2726                                  & 5.5           &                                        & 0.161                  & 0.117                  & 56.2           & 0.3793      \\
    15  & 0.05 & 2703                                  & 13.2          &                                        & 0.156                  & 0.095                  & 45.7           & 0.4769      \\
    15  & 0.10 & 2504                                  & 22.5          &                                        & 0.167                  & 0.072                  & 38.2           & 0.5941      \\
    15  & 0.15 & 2045                                  & 30.5          &                                        & 0.155                  & 0.044                  & 30.6           & 0.7781      \\
    \bottomrule
  \end{tabular}
  \datasource{Same source as Table~\ref{tab:x2-tracking}; columns
    \texttt{acid\_txn\_admitted\_count},
    \texttt{acid\_dirty\_block\_time\_mean\_ms},
    \texttt{acid\_fraction\_time}, \texttt{acid\_fraction\_completed} and
    \texttt{txn\_latency\_mean\_ms}. Every engaged cell of the
    sweep is listed; the five degenerate cells have no ACID admissions and are
    omitted.}
\end{table}

\paragraph{Reading the table.} Rows are grouped by reconciler count, and
within a group the budget increases downwards. $A$ is how many root
transactions were admitted in ACID mode. \emph{Mean wait} is the average
length of one wait on a repair. A transaction is parked in one of two places:
at admission, before it takes any locks, or mid-flight, when a bookkeeper
transaction it spawned finds a resource that has turned dirty since
(Section~\ref{sec:prelim-cprs}). Both are timed from parking until the
reconcilers release the resource, and the average is taken per wait rather
than per transaction, since one transaction can be made to wait more than
once. $\phi_{\mathrm{time}}$ and $\phi_{\mathrm{comp}}$ are the two ACID
shares of Section~\ref{sec:exp-metrics} --- of the run's time and of its
finished transactions --- and they are not interchangeable: the first says how
long the system stayed in the slow mode, the second how much of the finished
work was done there. \emph{Mean lat.} is the average end-to-end time from
admission to commit, over all completed transactions of either mode.

\paragraph{What the table shows.} The left half is the whole of the first
claim. At $R = 10$ the mean wait stretches from 7.9\,s to 59.0\,s as the budget
is raised, so ACID work gets \emph{worse}, not cheaper --- and there is
steadily less of it, with admissions falling from 2805 to 1320 over the same
rows. Every reconciler count behaves the same way. The right half says where
the extra throughput comes from anyway: $\phi_{\mathrm{comp}}$ falls monotonically at every reconciler
count --- 0.472 to 0.021 at $R = 1$, 0.108 to 0.024 at $R = 10$ --- so the
finished work is increasingly BASE work. Because an ACID transaction occupies
far more time than a BASE one, that shift pulls the mean latency down (50.6\,s
to 26.3\,s at $R = 10$), and under \texttt{EQUILIBRIUM} the in-flight
population is held constant (Section~\ref{sec:exp-workload}), so a transaction
that finishes sooner frees its slot sooner and throughput rises.

$\phi_{\mathrm{time}}$ is what makes the exchange a poor one. Where
$\phi_{\mathrm{comp}}$ collapses, $\phi_{\mathrm{time}}$ barely moves: at
$R = 10$ it stays between 0.24 and 0.29 across the whole budget range, and at
$R = 15$ between 0.155 and 0.167. The system is in ACID about as often as
before; it simply gets less done while it is there. The operator pays real
inconsistency and receives a reallocation of work from ACID to BASE inside an
ACID occupancy that stays where it was --- which is why the gain is real but
small, and why it cannot keep going once $\phi_{\mathrm{comp}}$ is near zero.

$R = 1$ is the case where even that fails. There $\phi_{\mathrm{time}}$ does
fall usefully, from 0.837 to 0.665, but repair is too slow for BASE to take up
the slack: ACID completions collapse from 2138 to 116 while BASE completions
rise only from 2395 to 5298. That is why a larger budget buys so little
throughput in general, and nothing at all at $R = 1$.

\section{X3: Removing the feedback signal}\label{sec:x3-random-baseline}

X3 is a single-variable removal. The random oracle keeps everything the
adaptive one has --- the arbiter, per-transaction tagging, the two
dispatchers --- but decides ACID or BASE by a coin flip with fixed
probability $p$ instead of by reading the dirty fraction. If the X2 results
survive that removal, they did not come from the feedback loop. Removing
that one input has consequences beyond the decision itself: without a
signal to hold it in ACID until the dirty set has drained, the random
oracle also stops collecting its ACID time in long excursions. That is a
downstream effect of the removal, not a second thing removed.

One property of the workload should be stated before the results rather
than after them. Every run in this thesis is executed under
\texttt{EQUILIBRIUM}, a stationary load (Section~\ref{sec:exp-workload}).
Feedback control earns its keep when the operating point moves, and a
constant one cannot reward it for tracking. The throughput comparison
below is therefore run in the regime least favourable to the adaptive
oracle, and a null result there is the expected one rather than a
surprising one. What that buys is a clean attribution: whatever advantage
adaptive control turns out to have under moving load
(Section~\ref{sec:future-work}) cannot be credited to hybrid execution in
general, because hybrid execution on its own is measured here and does not
supply it.

The sweep runs $p \in \{0.05, 0.10, 0.15, 0.25, 0.50\}$ against
$R \in \{1, 5, 10, 15\}$, five seeds per cell, 100 runs, everything else at
the baseline of Table~\ref{tab:exp-baseline}. The grid is concentrated
below 0.25 rather than centred on 0.5 because X1 shows the adaptive oracle
finishing between 7\% and 10\% of its work in ACID
(Table~\ref{tab:x1-summary}), and a grid centred on 0.5 would never bracket
that. Pure BASE and pure ACID are the $p = 0$ and $p = 1$ anchors, and are
already measured in X1. The adaptive comparison points are the 19 engaged
cells of X2.

\subsection{Results of the probability sweep}\label{sec:x3-results}

\begin{table}[htbp]
  \centering
  \caption{X3: results by ACID probability $p$ and reconciler count $R$.
    Columns are those of Table~\ref{tab:x1-summary}: $\phi_{\mathrm{time}}$
    is the share of the run spent in ACID mode, $\phi_{\mathrm{comp}}$ the
    share of completed transactions that ran under ACID, and
    \emph{Switches} how often the oracle changed the mode.}
  \label{tab:x3-summary}
  \small
  \begin{tabular}{r r r r r r r r}
    \toprule
    $R$ & $p$  & Thr.\ (tps) & $\phi_{\mathrm{time}}$ & $\phi_{\mathrm{comp}}$ & Switches & $p_{50}$ (\%) & $p_{95}$ (\%) \\
    \midrule
    1   & 0.05 & 0.0904      & 0.5988                 & 0.0445                 & 438.0    & 9.42          & 16.73         \\
        & 0.10 & 0.0760      & 0.6621                 & 0.0934                 & 698.6    & 6.96          & 12.22         \\
        & 0.15 & 0.0726      & 0.6789                 & 0.1397                 & 918.4    & 5.67          & 10.14         \\
        & 0.25 & 0.0733      & 0.7182                 & 0.2400                 & 1374.0   & 4.00          & 7.13          \\
        & 0.50 & 0.0831      & 0.7890                 & 0.4974                 & 2090.8   & 1.02          & 2.98          \\
    \midrule
    5   & 0.05 & 0.3918      & 0.1288                 & 0.0484                 & 1857.4   & 8.60          & 14.90         \\
        & 0.10 & 0.3317      & 0.2142                 & 0.0987                 & 2991.8   & 5.37          & 10.50         \\
        & 0.15 & 0.2731      & 0.3045                 & 0.1471                 & 3444.2   & 1.94          & 7.80          \\
        & 0.25 & 0.1711      & 0.4980                 & 0.2456                 & 3198.4   & 0.05          & 1.95          \\
        & 0.50 & 0.0849      & 0.7873                 & 0.4971                 & 2135.8   & 0.00          & 0.18          \\
    \midrule
    10  & 0.05 & 0.6684      & 0.0622                 & 0.0483                 & 3123.8   & 6.73          & 12.86         \\
        & 0.10 & 0.4254      & 0.1549                 & 0.0988                 & 3828.0   & 0.36          & 6.42          \\
        & 0.15 & 0.2839      & 0.2829                 & 0.1475                 & 3590.8   & 0.05          & 1.53          \\
        & 0.25 & 0.1707      & 0.4858                 & 0.2451                 & 3187.8   & 0.00          & 0.29          \\
        & 0.50 & 0.0850      & 0.7761                 & 0.4969                 & 2137.6   & 0.00          & 0.09          \\
    \midrule
    15  & 0.05 & 0.7920      & 0.0548                 & 0.0487                 & 3706.2   & 2.22          & 12.73         \\
        & 0.10 & 0.4416      & 0.1143                 & 0.0992                 & 3973.4   & 0.09          & 2.50          \\
        & 0.15 & 0.2942      & 0.1961                 & 0.1480                 & 3728.6   & 0.02          & 0.43          \\
        & 0.25 & 0.1750      & 0.3594                 & 0.2454                 & 3269.8   & 0.00          & 0.22          \\
        & 0.50 & 0.0876      & 0.6601                 & 0.4972                 & 2202.0   & 0.00          & 0.07          \\
    \bottomrule
  \end{tabular}
  \datasource{\texttt{random\_p\_sweep/random\_p\_sweep\_seed\{0..4\}.csv},
    100 runs, mean over 5 seeds. Seed CV on $p_{95}$: 5.52\% median,
    19.02\% max.}
\end{table}

\paragraph{The random oracle controls what it reads.} Compare the
$\phi_{\mathrm{comp}}$ entries for one $p$ across the four blocks and they
barely move: at $p = 0.05$ the realised ACID share runs 0.0445 to 0.0487, a
spread of 9\%, and at $p = 0.50$ it runs 0.4969 to 0.4974, a spread of
0.1\%. Nothing about the system enters that number. X1 reports the same
thing at a single probability (O4); the sweep shows it holds at all five.

\paragraph{It does not control the dirty fraction.} The $p_{95}$ column of
the same rows moves by a factor of 4.9 at $p = 0.10$ (12.22\% down to
2.50\%) and a factor of 43.9 at $p = 0.50$ (2.98\% down to 0.07\%). The
knob is set once; what the operator gets from it depends on how many
reconcilers happen to be provisioned. This is the sense in which the
random oracle loses the point of a critical dirty fraction: $p$ is not
expressed in the units of the quantity it is meant to bound, and no
setting of it names a consistency level.

\paragraph{But it is still a knob.} Inconsistency falls monotonically in
$p$ at every $R$, and monotonically in $R$ at every $p$. Under
\texttt{EQUILIBRIUM} the in-flight population is held constant
(Section~\ref{sec:exp-workload}), so raising $p$ replaces fast BASE
transactions with slow ACID ones, the admission rate drops, and less
dirtying work enters the system. The random oracle therefore does control
inconsistency --- indirectly, through the ACID share, rather than by
aiming at a target. That is what makes the two oracles comparable at all,
and Section~\ref{sec:x3-frontier} uses it.

\begin{figure}[H]
  \centering
  \includegraphics[width=\textwidth]{x3_control.pdf}
  \caption{Measured $p_{95}$ dirty fraction against reconciler count, one
    line per knob setting, on a shared logarithmic scale. Each oracle holds
    fixed the quantity it reads: the adaptive lines are flat because the
    threshold is a target, and the random lines fan out because the
    probability is not.}
  \label{fig:x3-control}
\end{figure}

Figure~\ref{fig:x3-control} puts the two side by side, and the adaptive
panel is the mirror image of the random one. Across the same range of $R$,
the measured $p_{95}$ under the adaptive oracle stays within 0.74\,pp of
the configured $c$ (Section~\ref{sec:x2-tracking}) while its realised ACID
share is what moves --- $\phi_{\mathrm{comp}}$ spans 0.108 to 0.472 at
$c = 0.02$, a factor of 4.4. Each oracle pins the quantity it reads and
lets the other one float. RQ3 is asked in the units of the dirty fraction,
so which quantity gets pinned is the whole of the difference.

\paragraph{Two columns to read with care.} $\phi_{\mathrm{time}}$ means
here exactly what it means in X1 and X2. Both oracles accrue it with the
same routine, which attributes the interval between two consecutive
decisions to whichever mode was in force across it, so in both cases the
column is the share of simulated time the admission mode was ACID, and in
both cases dividing it by the number of ACID entries gives the mean length
of a stay. What differs is how the share is assembled. The random oracle
changes mode 438 to 3973 times per run against 36 to 509 for the adaptive
one, so it collects its ACID time in many short visits rather than a few
long ones: mean stays of 1.5 to 137\,s, against 61 to 1835\,s under the
adaptive oracle over the same $R$.

The column should not be read against $p$. A probability of 0.05 says that
one transaction in twenty is routed to ACID, not that one twentieth of the
clock is spent there; ACID work is slow, so the few transactions that get
it hold the system in ACID for a large share of the run. That is the same
gap between $\phi_{\mathrm{time}}$ and $\phi_{\mathrm{comp}}$ described in
Section~\ref{sec:exp-metrics}, and it is widest here, at 0.60 against
$\phi_{\mathrm{comp}} = 0.045$ for $p = 0.05$ at $R = 1$.

Throughput at $R = 1$ is also not monotone in $p$: it falls from 0.0904 to
0.0726 and then rises again to 0.0831, while at $R = 5$, 10 and 15 it falls
steadily as $p$ rises. The reason is that $R$ splits a fixed pool rather
than adding to it (Section~\ref{sec:x1-reconciler-sweep}). Every reconciler
is taken out of the 30 foreground processors, so $R = 1$ leaves 29 running
transactions and one repairing, and $R = 15$ leaves 15 of each. The two
ends of the range are therefore limited by opposite things.

At $R = 1$ there is ample capacity to admit and execute work and almost
none to clean up after it: the dirty set drains at $0.5R = 0.5$ resources
per second (Section~\ref{sec:recovery-mechanism}) against a standing
backlog of over a thousand. An ACID transaction that touches a dirty
resource waits, and that wait --- not the choice of mode --- sets the pace
of the run. Moving work between the modes cannot speed up something limited
by waiting on repair, which is why $p$ has so little purchase here. At
$R = 15$ the dirty set drains fifteen times faster and waits are short, but
half the pool is now repairing instead of executing, so what limits the run
is the capacity left to process transactions at all --- and the mode mix,
which sets how expensive each transaction is, is exactly what moves that.
Throughput therefore falls monotonically in $p$ at the top of the range and
not at the bottom. The same mechanism produces the adaptive throughput
valley of X1 (Section~\ref{sec:x1-engagement}), with $c$ rather than $p$
setting the mix.

\subsection{Throughput at matched inconsistency}\label{sec:x3-frontier}

RQ3 asks whether reading the metrics beats a coin flip \emph{on the same
  throughput/inconsistency trade-off}. The difficulty is that the two
oracles are configured in different terms. Setting $c = 0.10$ asks the
adaptive oracle to hold dirtiness near 10\%; setting $p = 0.10$ tells the
random oracle to send one transaction in ten to ACID. Those are not the
same instruction, and no pair of settings makes the two runs comparable.

What is comparable is what the runs produced. Every run reports both a
throughput and the dirtiness it actually reached, so the question can be
put on the output side: at the dirtiness this adaptive run reached, how
fast was the random oracle? Answering it needs a random run at that same
dirtiness, and the random arm was swept at only five probabilities. The
five seeds of each configuration are averaged first, so each setting
contributes one throughput and one dirtiness.

\paragraph{Locating each adaptive run among the random ones.} Take the
adaptive run at $R = 5$, $c = 0.10$, which ended up 10.22\% dirty at
0.345\,tps. The five random runs at $R = 5$ ended up 14.90, 10.50, 7.80,
1.95 and 0.18\% dirty. None of them is at 10.22\%, but two of them sit
either side of it --- the run at 10.50\% and the run at 7.80\% --- and the
speed at both is known. Table~\ref{tab:x3-brackets} does this for every
adaptive run: it records where the run landed, which two random runs it
fell between, and how close it came to the nearer of the two. That last
figure matters because a straight-line estimate is exact at its endpoints
and least reliable midway between them: an adaptive run that lands almost
on top of a measured random run barely needs the estimate at all. It is
reported as a ratio --- the adaptive run at 5.43\% sits between runs at
2.98 and 7.13\%, and $7.13 / 5.43 = 1.31$, so it is 31\% away from the
nearer of the two. A value of 1.00 means the two coincide.

\begin{table}[htbp]
  \centering
  \caption{Where each adaptive run falls among the random runs at the same
    reconciler count. \emph{Between} gives the dirtiness of the two random
    runs on either side; the full set of random runs is the $p_{95}$ column
    of Table~\ref{tab:x3-summary}. \emph{To nearest} is how far the adaptive
    run sits from the closer of those two, as a ratio --- 1.00 means it
    landed on top of a measured run. A run with no random run on one side
    cannot be compared and is dropped rather than extrapolated.}
  \label{tab:x3-brackets}
  \small
  \begin{tabular}{r r r l c}
    \toprule
    $R$ & $c$  & adaptive $p_{95}$ (\%) & Between (\%)                         & To nearest   \\
    \midrule
    1   & 0.02 & 2.52                   & \emph{cleaner than every random run} & ---          \\
    1   & 0.05 & 5.43                   & 2.98 -- 7.13                         & 1.31$\times$ \\
    1   & 0.10 & 10.12                  & 7.13 -- 10.14                        & 1.00$\times$ \\
    1   & 0.15 & 14.76                  & 12.22 -- 16.73                       & 1.13$\times$ \\
    1   & 0.20 & 19.52                  & \emph{dirtier than every random run} & ---          \\
    \midrule
    5   & 0.02 & 2.32                   & 1.95 -- 7.80                         & 1.19$\times$ \\
    5   & 0.05 & 5.36                   & 1.95 -- 7.80                         & 1.46$\times$ \\
    5   & 0.10 & 10.22                  & 7.80 -- 10.50                        & 1.03$\times$ \\
    5   & 0.15 & 14.79                  & 10.50 -- 14.90                       & 1.01$\times$ \\
    5   & 0.20 & 19.54                  & \emph{dirtier than every random run} & ---          \\
    \midrule
    10  & 0.02 & 2.10                   & 1.53 -- 6.42                         & 1.37$\times$ \\
    10  & 0.05 & 5.14                   & 1.53 -- 6.42                         & 1.25$\times$ \\
    10  & 0.10 & 9.98                   & 6.42 -- 12.86                        & 1.29$\times$ \\
    10  & 0.15 & 14.56                  & \emph{dirtier than every random run} & ---          \\
    10  & 0.20 & 19.44                  & \emph{dirtier than every random run} & ---          \\
    \midrule
    15  & 0.02 & 1.84                   & 0.43 -- 2.50                         & 1.36$\times$ \\
    15  & 0.05 & 4.71                   & 2.50 -- 12.73                        & 1.88$\times$ \\
    15  & 0.10 & 9.36                   & 2.50 -- 12.73                        & 1.36$\times$ \\
    15  & 0.15 & 14.28                  & \emph{dirtier than every random run} & ---          \\
    \bottomrule
  \end{tabular}
  \datasource{\texttt{random\_p\_sweep/random\_p\_sweep\_seed\{0..4\}.csv}
    and \texttt{critical\_sweep\_seed\{0..4\}.csv}, means over 5 seeds.
    $R = 15$, $c = 0.20$ is degenerate and omitted
    (Section~\ref{sec:x2-critical-sweep}).}
\end{table}

Thirteen of the 19 engaged adaptive runs fall between two random runs and
can be compared. The remaining six have nothing to sit between --- five
are dirtier than every random run at that reconciler count, and one, at
$R = 1$, $c = 0.02$, is cleaner than every one --- and are dropped rather
than extrapolated.

\paragraph{Reading the speed in between.} For the thirteen that can be
compared, the random oracle's speed at the adaptive run's dirtiness is
estimated between the two runs on either side. Continuing the example
above: the random run at 10.50\% dirty managed 0.332\,tps and the one at
7.80\% managed 0.273\,tps, and 10.22 lies nine tenths of the way from 7.80
up to 10.50, so the estimate is nine tenths of the way from 0.273 up to
0.332 --- 0.326\,tps. The adaptive run's 0.345 against that is a ratio of
1.06, the same throughput to within a few percent.
Table~\ref{tab:x3-matched} is that calculation for all thirteen.

How far such an estimate can be trusted depends on the \emph{To nearest}
column of Table~\ref{tab:x3-brackets}. In the example the adaptive run at
10.22\% sat 3\% below the measured run at 10.50\%, so the estimate hardly
moves off a measured value and the comparison is close to a direct one.
Where the nearest measured run is half again as far away, the estimate is
carrying the answer, and the row settles nothing.

\begin{table}[htbp]
  \centering
  \caption{Throughput of the two oracles at the same measured
    inconsistency, for the thirteen adaptive runs that fall between two
    random runs. \emph{Ratio} is adaptive divided by random. Rows are
    ordered by \emph{To nearest} (Table~\ref{tab:x3-brackets}), so the
    comparisons that are closest to a straight measurement come first; the
    rule marks where the adaptive run stops sitting within 15\% of a
    measured one.}
  \label{tab:x3-matched}
  \small
  \begin{tabular}{r r r r r c c}
    \toprule
        &      & matched       & \multicolumn{2}{c}{Thr.\ (tps)} &        &                      \\
    \cmidrule(lr){4-5}
    $R$ & $c$  & $p_{95}$ (\%) & adaptive                        & random & Ratio & To nearest   \\
    \midrule
    1   & 0.10 & 10.1          & 0.073                           & 0.073  & 1.00  & 1.00$\times$ \\
    5   & 0.15 & 14.8          & 0.397                           & 0.390  & 1.02  & 1.01$\times$ \\
    5   & 0.10 & 10.2          & 0.345                           & 0.326  & 1.06  & 1.03$\times$ \\
    1   & 0.15 & 14.8          & 0.084                           & 0.084  & 1.00  & 1.13$\times$ \\
    \midrule
    5   & 0.02 & 2.3           & 0.255                           & 0.178  & 1.44  & 1.19$\times$ \\
    10  & 0.05 & 5.1           & 0.486                           & 0.388  & 1.25  & 1.25$\times$ \\
    10  & 0.10 & 10.0          & 0.588                           & 0.560  & 1.05  & 1.29$\times$ \\
    1   & 0.05 & 5.4           & 0.072                           & 0.077  & 0.93  & 1.31$\times$ \\
    15  & 0.02 & 1.8           & 0.379                           & 0.395  & 0.96  & 1.36$\times$ \\
    15  & 0.10 & 9.4           & 0.594                           & 0.676  & 0.88  & 1.36$\times$ \\
    10  & 0.02 & 2.1           & 0.395                           & 0.300  & 1.32  & 1.37$\times$ \\
    5   & 0.05 & 5.4           & 0.303                           & 0.231  & 1.31  & 1.46$\times$ \\
    15  & 0.05 & 4.7           & 0.477                           & 0.517  & 0.92  & 1.88$\times$ \\
    \bottomrule
  \end{tabular}
  \datasource{Same two sources as Table~\ref{tab:x3-brackets}, means over
    5 seeds.}
\end{table}

\paragraph{Where the comparison is nearly a measurement, the two are
  equal.} Four adaptive runs land within 15\% of a measured random run,
and all four give the same throughput as the coin flip: 1.00, 1.02, 1.06
and 1.00. These are the rows where almost nothing rests on the estimate,
and they are unanimous. All four are at $c \ge 0.10$, because that is
where the random sweep happens to have runs nearby.

\paragraph{The other rows point in no direction.} The remaining nine
adaptive runs land too far from any measured random run for a single row
to be worth much on its own. Taken together they say nothing either: five
come out above parity and four below it, the median is 1.05, and no
setting of $c$ or of the reconciler count turns one oracle into the
reliable winner. Whatever is moving those numbers, it is not a benefit
belonging to either oracle. A finer grid of probabilities would let these
rows be read one by one rather than only in bulk, and is left to future
work (Section~\ref{sec:future-work}).

\paragraph{The answer to RQ3.} Reading the mode from the system's own
metrics does not beat a coin flip in any way this experiment can
establish. Where the two oracles can be compared against a nearby
measurement they are equal, and everywhere else the sweep is too coarse to
say. A finer grid of probabilities would settle the remaining rows one way
or the other, and is left to future work.

None of this is the setting in which feedback would be expected to prove
itself. Every run in this thesis is executed under a stationary load, and
a controller that adapts has nothing to adapt to when the operating point
never moves. Testing whether the feedback loop earns its keep needs a
workload that changes while the run is in progress --- the terminal-driven
mode the simulator already supports, left to future work
(Section~\ref{sec:future-work}). Until that is run, the case for reading
the signal is the one made in Section~\ref{sec:x3-results}, which does not
rest on throughput at all.

\subsection{Findings}\label{sec:x3-findings}

\begin{description}
  \item[O9] \textbf{Each oracle controls the quantity it reads, and only
          that one.} The random oracle delivers the ACID share it was told to
        deliver, whatever the system is doing, but the inconsistency that
        results swings widely with the reconciler count. The adaptive oracle
        is the mirror image: it holds the inconsistency where it was asked to
        and lets the ACID share fall where it may. From this we can
        understand that neither oracle controls both quantities, and that
        choosing between them is a choice of which one the operator wants to
        be able to name.

  \item[O10] \textbf{Neither oracle is measurably faster than the other
          under a stationary load.} Where an adaptive run lands close enough to
        a measured random run for the comparison to be nearly direct, the two
        turn in the same throughput; where the data hints otherwise, the
        random sweep is too coarse to support the hint. From this we can understand
        that the feedback loop costs nothing in throughput, and that this
        workload cannot show whether it is worth anything either: with the
        operating point held constant there is nothing for an adaptive
        controller to adapt to. A varying workload is the test that would
        settle it.

\end{description}

\section{X4: Semantic overlap degree sweep}\label{sec:x4-overlap-sweep}

\emph{Semantic overlap} is the parameter Zacouris and
Acosta~\cite{zacouris2025m2ts} use to describe how tightly the models of a
multi-model system are tied to one another
(Section~\ref{sec:prelim-cprs}). Two models overlap semantically when they
hold resources standing for the same underlying fact, and a consistency
preservation rule ties those resources together. The overlap degree
$\sigma$ (\texttt{semanticOverlap}) is the fraction of the resources in
the system that carry such a cross-model rule. The simulator draws that
set once at start-up, from a seeded draw, and keeps it fixed for the whole
run: a write to a resource in the set spawns a bookkeeper transaction
against the resource it overlaps in another model, and leaves that
resource dirty until a reconciler clears it, while a write to any other
resource stays inside its own model and dirties nothing. So $\sigma$ is a
property of the schema, not of the load --- it says how much of the rest
of the system one write disturbs.

The workload of this sweep is the one used in X1 to X3, unchanged: the
same arrival rate, the same transaction size $L = 50$, the same resource
count. Only $\sigma$ moves. That is what makes the sweep worth running on
its own: it changes the rate at which the dirty set fills --- the plant
the oracle regulates --- without changing how much work is offered to the
system. The budget of X2 was established at a single value of $\sigma$,
and an operator needs to know over what range of overlap it still
holds.

The sweep covers $\sigma \in \{0.00, 0.05, 0.10, 0.20, 0.30, 0.50,
  1.00\}$ against $R \in \{1, 5, 10, 15\}$ over five seeds, 140 runs, under
\texttt{HYBRID\_ADAPTIVE} at the X1 baseline thresholds $c = 0.10$,
$r = 0.05$. The $\sigma = 0.30$ block is the shared baseline of every
other experiment and reproduces the X1 rows exactly, which is the
cross-check that the two sweeps agree where they overlap.

\begin{table}[htbp]
  \centering
  \caption{X4: dirty-fraction percentiles, oracle behaviour, throughput and
    abort rate against the semantic overlap degree $\sigma$, by reconciler
    count $R$. ``deg.'' marks a cell in which the oracle never switches, so
    the row is pure BASE under a different label. Abort rate is
    \texttt{aborted\_total} over admitted-plus-aborted.}
  \label{tab:x4-overlap}
  \small
  \begin{tabular}{r r r r r r r r r}
    \toprule
    $\sigma$ & $R$ & $p_{50}$ (\%) & $p_{95}$ (\%) & Switches & $\phi_{\mathrm{time}}$ & $\phi_{\mathrm{comp}}$ & Thr.\ (tps) & Aborts \\
    \midrule
    0.00     & 1   & 0.00          & 0.00          & deg.     & 0.0000                 & 0.0000                 & 3.8685      & 0.000  \\
    0.00     & 5   & 0.00          & 0.00          & deg.     & 0.0000                 & 0.0000                 & 3.3344      & 0.000  \\
    0.00     & 10  & 0.00          & 0.00          & deg.     & 0.0000                 & 0.0000                 & 2.6681      & 0.000  \\
    0.00     & 15  & 0.00          & 0.00          & deg.     & 0.0000                 & 0.0000                 & 2.0013      & 0.000  \\
    \midrule
    0.05     & 1   & 4.50          & 4.62          & deg.     & 0.0000                 & 0.0000                 & 3.8685      & 0.000  \\
    0.05     & 5   & 2.14          & 2.44          & deg.     & 0.0000                 & 0.0000                 & 3.3344      & 0.000  \\
    0.05     & 10  & 0.08          & 0.22          & deg.     & 0.0000                 & 0.0000                 & 2.6681      & 0.000  \\
    0.05     & 15  & 0.04          & 0.15          & deg.     & 0.0000                 & 0.0000                 & 2.0013      & 0.000  \\
    \midrule
    0.10     & 1   & 9.50          & 9.62          & deg.     & 0.0000                 & 0.0000                 & 3.8685      & 0.000  \\
    0.10     & 5   & 7.16          & 7.47          & deg.     & 0.0000                 & 0.0000                 & 3.3344      & 0.000  \\
    0.10     & 10  & 3.20          & 3.66          & deg.     & 0.0000                 & 0.0000                 & 2.6681      & 0.000  \\
    0.10     & 15  & 0.14          & 0.35          & deg.     & 0.0000                 & 0.0000                 & 2.0013      & 0.000  \\
    \midrule
    0.20     & 1   & 6.54          & 9.90          & 82       & 0.6746                 & 0.0601                 & 0.1425      & 0.025  \\
    0.20     & 5   & 5.02          & 9.92          & 218      & 0.3856                 & 0.0670                 & 0.6495      & 0.036  \\
    0.20     & 10  & 2.02          & 9.60          & 195      & 0.1951                 & 0.0526                 & 0.9897      & 0.030  \\
    0.20     & 15  & 7.54          & 8.22          & deg.     & 0.0000                 & 0.0000                 & 2.0013      & 0.000  \\
    \midrule
    0.30     & 1   & 6.80          & 10.12         & 80       & 0.7776                 & 0.0918                 & 0.0726      & 0.070  \\
    0.30     & 5   & 5.67          & 10.22         & 234      & 0.5068                 & 0.0966                 & 0.3453      & 0.074  \\
    0.30     & 10  & 3.22          & 9.98          & 224      & 0.2886                 & 0.0694                 & 0.5883      & 0.074  \\
    0.30     & 15  & 0.43          & 9.36          & 162      & 0.1667                 & 0.0716                 & 0.5941      & 0.052  \\
    \midrule
    0.50     & 1   & 7.22          & 10.26         & 68       & 0.8767                 & 0.1384                 & 0.0299      & 0.154  \\
    0.50     & 5   & 6.42          & 10.40         & 213      & 0.6423                 & 0.1314                 & 0.1421      & 0.132  \\
    0.50     & 10  & 4.95          & 10.21         & 240      & 0.4352                 & 0.0945                 & 0.2581      & 0.146  \\
    0.50     & 15  & 2.40          & 9.95          & 192      & 0.3023                 & 0.0773                 & 0.3127      & 0.137  \\
    \midrule
    1.00     & 1   & 9.50          & 22.48         & 19       & 0.9784                 & 0.2900                 & 0.0069      & 0.252  \\
    1.00     & 5   & 9.49          & 12.14         & 78       & 0.8836                 & 0.2351                 & 0.0274      & 0.337  \\
    1.00     & 10  & 7.46          & 11.34         & 133      & 0.7830                 & 0.1464                 & 0.0446      & 0.338  \\
    1.00     & 15  & 5.73          & 10.79         & 159      & 0.6802                 & 0.1042                 & 0.0607      & 0.282  \\
    \bottomrule
  \end{tabular}
  \datasource{\texttt{semantic\_overlap\_sweep/semantic\_overlap\_sweep\_seed\{0..4\}.csv},
    140 runs, mean over 5 seeds. Seed CV on $p_{95}$: 0.41\% median, 7.21\%
    maximum.}
\end{table}

\subsection{Reading the sweep}\label{sec:x4-engagement}\label{sec:x4-invariance}\label{sec:x4-saturation}

Table~\ref{tab:x4-overlap} is read downwards in blocks. Each block holds
$\sigma$ fixed and steps through the reconciler count, so a column read
across blocks shows what overlap does at fixed repair capacity, and a
column read inside one block shows what repair capacity does at fixed
overlap. Three things have to be separated when reading it. Whether the
oracle fires at all is the switch count and the ``deg.'' marker: a
degenerate cell is pure BASE under a different label, and averaging over
it would mix two systems. Whether the oracle meets its budget is
$p_{95}$ against the configured $c = 10\%$. What that costs is the
throughput and abort columns. The three move independently, and the
sweep is informative because each is decided by a different part of the
system.

\paragraph{Engagement has a threshold in $\sigma$.} Thirteen of the 28
cells are degenerate, and they are not scattered: every cell at
$\sigma \le 0.10$ degenerates at every $R$, plus $\sigma = 0.20$ at
$R = 15$. This is the degeneracy rule of finding O6
(Section~\ref{sec:x2-degenerate}) again, read along a different axis.
Coupling density sets the BASE steady-state dirty fraction, and the
oracle engages only where that steady state reaches $c$. At
$\sigma = 0.10$, BASE alone runs at $p_{95} = 9.62\%$ at $R = 1$, already
below $c = 0.10$ and falling to 0.35\% at $R = 15$; at $\sigma = 0.20$ it
clears $c$ at $R \le 10$ and not at $R = 15$. The boundary is sharp:
between $\sigma = 0.10$ and $\sigma = 0.20$ the sweep goes from zero
switches in twenty runs to 82--218 switches per run. The practical form
of this is that a hybrid deployment on a schema whose models barely
overlap is pure BASE regardless of configuration.

\paragraph{Where it engages, the budget does not move with overlap.}
Once the oracle is firing, more overlap does not make the system more
inconsistent. It makes the oracle work harder, and the tail comes out
where it was configured. The eleven engaged cells at $\sigma \in \{0.20,
  0.30, 0.50\}$ all land within 0.64\,pp of the configured $c = 10\%$, and
neither $\sigma$ nor $R$ moves the result in any direction inside that
band. This is the strongest form the X2 claim takes anywhere in this
chapter: the threshold is a budget on the realised tail, and how much the
models overlap does not enter it.

\paragraph{What $\sigma$ does control is the price.} What rises with
overlap is not the inconsistency but the cost of keeping it inside the
budget. At $R = 1$ throughput falls steadily --- at $\sigma = 0.20$:
0.14\,tps; at $\sigma = 0.30$: 0.073\,tps; at $\sigma = 0.50$:
0.030\,tps; at $\sigma = 1.00$: 0.0069\,tps --- a factor of 21 across the
engaged range, and a factor of 560 against the 3.87\,tps the same $R$
delivers where the oracle never fires. Aborts climb with it, from 0.025
to 0.25, while the tail itself moves by 0.36\,pp. Holding the budget
fixed as overlap rises is paid for entirely in throughput and aborts.

\paragraph{One cell breaks the invariance, and it is the interesting
  one.} At $\sigma = 1.00$, $R = 1$ the tail reaches 22.48\%, twice the
budget and the only such miss in this chapter. It is not a control
failure, and the entry measurements say so: the oracle switches to ACID
at a dirty fraction of 10.5\% on average, on the line it was given, and
everything above that arrives after the brake is already on. A switch
changes the mode of transactions admitted after it, not of the thirty
already in flight, and at $\sigma = 1.00$ what those thirty still have to
deposit is enormous. Over the run, 244 BASE root transactions produce
25\,397 dirty-mark events --- about 104 resources from a single
transaction, some 2\% of the whole system --- so thirty of them in flight
carry on the order of 3000 marks already committed to happen, against a
threshold that sits at 500 resources. The threshold is crossed by a
fraction of that output and the rest lands regardless. What decides the
tail is then how long the overshoot takes to unwind: the single
reconciler clears 25\,000 resources over the run, exactly its
$0.5$\,s$^{-1}$ flat out from start to finish, and the mean excursion
lasts 5239\,s. Reconcilers do not prevent the overshoot --- the peak is
still 18.5\% at $R = 15$ --- but they drain it fifteen times faster, and
$p_{95}$ comes back to 10.79\%. The limit is capacity, not the threshold.

\subsection{Findings}\label{sec:x4-findings}

\begin{description}
  \item[O11] \textbf{How much the models overlap decides whether the
          oracle has anything to do, not how well it does it.} Below a threshold in $\sigma$ the run is
        pure BASE at every reconciler count; above it, the tail lands within
        0.64\,pp of the configured budget whatever the overlap. From this we
        can understand that $\sigma$ is not a knob that trades consistency
        against throughput: what it changes is the price of the budget, which
        rises by a factor of 21 in throughput across the engaged range. An
        operator can carry one value of $c$ across schemas that overlap to
        different degrees, but not one throughput expectation.

  \item[O12] \textbf{The budget holds only while reconciliation keeps up
          with the dirtying it has to repair.} At $\sigma = 1.00$, $R = 1$ the
        oracle runs almost entirely in ACID and still overshoots $c$ by a factor
        of two, and
        the miss disappears as soon as $R$ rises. From this we can understand
        that the threshold is a conditional guarantee: once repair is slower
        than the dirtying, the oracle has already spent everything the policy
        allows, and no setting of $c$ or $r$ recovers it. The remedy there is
        provisioning, not tuning.
\end{description}

\section{X5: Resources per transaction sweep}\label{sec:x5-resources-sweep}

$L$ (\texttt{maxTransactionResources}) is the largest number of resources
a single transaction may touch. Each transaction draws its size uniformly
from $1$ to $L$, draws a model count uniformly from one to four
(\texttt{maxTransactionModels}, a configured baseline rather than a
property of the simulator, held at 4 by every experiment in this chapter),
and splits its resources over that many of the ten models, taking a
contiguous block inside each. That cap says how many models a transaction
reaches into, not how much of them it holds: at the baseline $L = 50$ the
block taken from one model averages ten of the 500 resources that model
owns. Nothing else moves with $L$. The schema is
the one every other experiment uses, ten models and 5000 resources at
$\sigma = 0.30$, so how far a single write propagates is unchanged; the
workload is the one every other experiment uses, \texttt{EQUILIBRIUM} with
30 transactions in flight, so the amount of work offered to the system is
unchanged too. What $L$ decides is only how deep one transaction reaches
into that fixed schema.

An example fixes the range. At the baseline $L = 50$ a transaction touches
25.5 resources on average out of the 5000 in the system, each under its
own lock until it commits, so the 30 transactions in flight hold some 750
between them, 15\% of the data set. At $L = 400$ the same draw averages
200.5, and the 30 in flight demand around 6000 --- more claims than the
system has locks to give. Nothing about the schema or about the arrival of
work differs between those two settings; each transaction simply reaches
deeper.

That makes this sweep a check rather than a new question. X2 established
the budget at one operating point and X4 showed it survives a change of
schema; here the axis is how deep a transaction goes, and the two things
worth measuring are whether the tail still lands at $c$ and, if it does,
which measured quantity absorbed the change instead. The sweep covers
$L \in \{5, 10, 25, 50, 100, 200, 400\}$ against $R \in \{1, 5, 10, 15\}$
over five seeds, 140 runs, under \texttt{HYBRID\_ADAPTIVE} at $c = 0.10$,
$r = 0.05$ and $\sigma = 0.30$. $L = 50$ is the baseline shared with every
other experiment.

Table~\ref{tab:x5-resources} is arranged for that pair of questions: one
block per reconciler count, rows ascending in $L$, so each block reads as
one repair capacity taken from small transactions to large ones. The
switch count says whether the oracle is engaged at all, $p_{95}$ against
the configured $c = 10\%$ is the quantity under control, $p_{99}$ says how
far the excursions above it reach, and the last three columns are what
holding the budget costs. Two of those need a word before the table.
\emph{Lat.} is the mean time from creation to commit
(Section~\ref{sec:exp-metrics}), which is what grows when a transaction
has more resources to process. \emph{Aborts} is \texttt{aborted\_total}
over admitted-plus-aborted, the share of started transactions that did not
finish.

\begin{table}[htbp]
  \centering
  \caption{X5: what moves and what does not as the resources per
    transaction $L$ grow, by reconciler count $R$.}
  \label{tab:x5-resources}
  \small
  \setlength{\tabcolsep}{4.5pt}
  \renewcommand{\arraystretch}{0.92}
  \begin{tabular}{r r r r r r r r r}
    \toprule
    $L$ & $R$ & $p_{50}$ (\%) & $p_{95}$ (\%) & $p_{99}$ (\%) & Switches & Thr.\ (tps) & Lat.\ (s) & Aborts \\
    \midrule
    5   & 1  & 7.35 & 9.74  & 9.98  & 111 & 0.6377 & 18.0  & 0.004 \\
    10  & 1  & 7.34 & 9.80  & 10.02 & 105 & 0.3458 & 40.5  & 0.012 \\
    25  & 1  & 7.12 & 9.96  & 10.27 & 92  & 0.1450 & 127.8 & 0.046 \\
    50  & 1  & 6.80 & 10.12 & 10.68 & 80  & 0.0726 & 261.8 & 0.070 \\
    100 & 1  & 6.46 & 10.29 & 11.92 & 66  & 0.0358 & 497.1 & 0.106 \\
    200 & 1  & 6.22 & 10.78 & 17.38 & 54  & 0.0177 & 820.2 & 0.283 \\
    400 & 1  & 5.91 & 11.87 & 21.46 & 43  & 0.0091 & 1221.3 & 0.522 \\
    \midrule
    5   & 5  & 7.08 & 9.77  & 10.02 & 457 & 3.0352 & 4.2   & 0.001 \\
    10  & 5  & 6.95 & 9.84  & 10.13 & 430 & 1.6682 & 9.2   & 0.006 \\
    25  & 5  & 6.44 & 10.07 & 10.46 & 325 & 0.6994 & 28.6  & 0.027 \\
    50  & 5  & 5.67 & 10.22 & 10.82 & 234 & 0.3453 & 67.2  & 0.074 \\
    100 & 5  & 4.00 & 10.19 & 11.03 & 147 & 0.1548 & 156.7 & 0.208 \\
    200 & 5  & 2.29 & 10.02 & 11.38 & 86  & 0.0639 & 334.0 & 0.403 \\
    400 & 5  & 0.94 & 9.48  & 11.95 & 45  & 0.0254 & 538.4 & 0.646 \\
    \midrule
    5   & 10 & 6.01 & 9.67  & 10.00 & 615 & 5.4951 & 3.1   & 0.000 \\
    10  & 10 & 5.29 & 9.72  & 10.13 & 522 & 2.9549 & 6.3   & 0.001 \\
    25  & 10 & 4.03 & 9.86  & 10.46 & 319 & 1.1752 & 16.7  & 0.016 \\
    50  & 10 & 3.22 & 9.98  & 10.82 & 224 & 0.5883 & 37.4  & 0.074 \\
    100 & 10 & 1.23 & 9.77  & 10.88 & 134 & 0.2447 & 98.2  & 0.195 \\
    200 & 10 & 0.58 & 9.14  & 10.70 & 69  & 0.0975 & 224.8 & 0.359 \\
    400 & 10 & 0.42 & 8.44  & 10.42 & 39  & 0.0439 & 369.9 & 0.532 \\
    \midrule
    5   & 15 & 5.14 & 9.60  & 9.94  & 493 & 7.5079 & 2.7   & 0.000 \\
    10  & 15 & 3.42 & 9.52  & 10.00 & 451 & 3.6856 & 6.1   & 0.000 \\
    25  & 15 & 0.83 & 9.44  & 10.16 & 254 & 1.2084 & 18.4  & 0.006 \\
    50  & 15 & 0.43 & 9.36  & 10.37 & 162 & 0.5941 & 38.2  & 0.052 \\
    100 & 15 & 0.42 & 9.22  & 10.59 & 113 & 0.2937 & 81.7  & 0.172 \\
    200 & 15 & 0.32 & 8.64  & 10.08 & 58  & 0.1294 & 181.0 & 0.318 \\
    400 & 15 & 0.52 & 7.91  & 9.70  & 31  & 0.0676 & 262.8 & 0.426 \\
    \bottomrule
  \end{tabular}
  \datasource{\texttt{max\_txn\_resources\_sweep/max\_txn\_resources\_sweep\_seed\{0..4\}.csv},
    140 runs, mean over 5 seeds. Seed CV on $p_{95}$: 0.26\% median, 4.23\%
    maximum.}
\end{table}

\subsection{Reading the sweep}\label{sec:x5-invariance}\label{sec:x5-cost}

Every cell switches, between 31 and 615 times, so no cell is degenerate
and all 28 are comparable. Down each block the $p_{95}$ column hardly
moves. Over the 24 cells with $L \le 200$ it spans 8.64\% to 10.78\%
against the configured $c = 10\%$, and over $L \in [5, 100]$, where the
resources per transaction grow by a factor of 20, it spans 9.22\% to
10.29\% --- within 0.78\,pp of the configured value. So $L$, which decides how many
resources a single transaction can dirty and therefore how fast the dirty
set fills, does not move the level at which that set is held.

What moves is in the cost columns, and neither of the two is the oracle's
doing. A transaction takes proportionally longer the more resources it
touches, because the simulator charges processing time per resource: at
$R = 1$, at $L = 5$: 18.0\,s from creation to commit; at $L = 400$:
1221\,s. With 30 transactions held in flight throughout, ones that take
longer individually complete less often, and that is the whole of the
throughput column: at $R = 1$, at $L = 5$: 0.638 transactions per second;
at $L = 400$: 0.0091. The second cost is the abort rate, which rises with
$L$ in every block: at $R = 10$, at $L = 5$: 0.000; at $L = 50$: 0.074; at
$L = 200$: 0.359; at $L = 400$: 0.532. Larger transactions hold more locks
at once, so they collide more often and more of the work already done is
thrown away.

Reconciliation is the one place where this outcome is worth an argument,
because the obvious expectation runs the other way: a transaction that
touches more resources dirties more of them, and should therefore load the
reconciler more. Per transaction it does, exactly in proportion. Dividing the
\texttt{dirty\_mark\_events\_total} of each run by the transactions it
completed gives, at $R = 1$, at $L = 5$: 0.8 dirty marks per transaction;
at $L = 400$: 53.
What breaks the inference is that the sweep holds 30 transactions in
flight rather than a fixed arrival rate, so the larger ones run
proportionally slower and the marks arrive no faster: at $R = 1$, at
$L = 5$: 0.509 marks per second; at $L = 400$: 0.486. The served side
agrees, and at high $R$
it moves against the expectation rather than merely failing to confirm it:
the clear rate of
Section~\ref{sec:exp-metrics} against its ceiling of $0.5R$ gives, at
$R = 15$, at $L = 5$: 7.45 resources cleared per second of a possible 7.5,
a saturated pool; at $L = 400$: 4.43, three fifths of capacity. Larger transactions give the
reconciler less to do, not more, because the lock conflicts they cause
hold down the rate at which resources are touched at all.

Two smaller patterns come with it. The median dirty fraction falls while
the tail stays where it was configured: at $R = 15$, at $L = 5$: 5.14\%;
at $L = 400$: 0.52\%, against a $p_{95}$ that stays between 7.91\% and
9.60\% throughout. With larger transactions the dirty set sits lower and
spikes harder, which is also what the falling switch counts show. And the
excursions above the budget grow only where repair is scarce: $p_{99}$ at
$R = 1$, at $L = 5$: 9.98\%; at $L = 400$: 21.46\%, while at $R \ge 10$ it
stays near 10\% at every $L$. That is O12 read on this axis
rather than on the overlap axis.

The $L = 400$ block sits outside the regime the rest of the sweep
measures. The 30 transactions in flight there demand about 6000 resources
against the 5000 locks the system has, so lock conflict rather than the
oracle sets the admission rate: the dirty set never gets the chance to
fill, and $p_{95}$
undershoots: at $R = 10$: 8.44\%; at $R = 15$: 7.91\%. The block is a
stress anchor, not a like-for-like scaling point, and is excluded from the
invariance claim above.

\subsection{Findings}\label{sec:x5-findings}

\begin{description}
  \item[O13] \textbf{How many resources a transaction touches does not move
          the tail.} As $L$ grows by a factor of 20 the measured $p_{95}$
        stays within 0.78\,pp of the configured $c$ at every reconciler
        count, with no degenerate cell. From this we can understand that O11
        was not an artefact of the axis it was measured on: the budget holds
        on both workload dimensions, the one that sets how far a write
        spreads and the one that sets how much a transaction touches.

  \item[O14] \textbf{What $L$ costs is latency and aborts, not reconciler
          load.} Larger transactions take proportionally longer to process and
        collide more often --- at $R = 1$, at $L = 5$: 18.0\,s per transaction
        and an abort rate of 0.004; at $L = 400$: 1221\,s and 0.52 --- while
        the dirty marks the reconciler has to clear arrive no faster, at
        $L = 5$: 0.509 per second; at $L = 400$: 0.486. From this we can understand that the
        two workload axes are priced differently: semantic overlap loads the
        reconciler and so raises the cost of the budget directly, while $L$ is
        paid for in the foreground, in waiting and in wasted work. An operator
        with a throughput problem should look at semantic overlap first, and
        at $L$ only if aborts are high.
\end{description}

\section{Ablation: the smoothing factor and its removal}\label{sec:alpha-ablation}

The five sweeps above evaluate one system. This section does not: it
reports the measurement that produced that system, by sweeping a parameter
the evaluated oracle no longer has. It is placed here rather than in the
approach chapter because the argument for the removal is empirical, and
placed after X5 rather than among the numbered sweeps because $\alpha$ is
not a knob any result in this chapter can be reproduced with.

The oracle evaluated in this thesis compares its thresholds against the
dirty fraction measured at the instant of the decision. An earlier version
of the implementation did not: it compared them against an exponentially
weighted moving average of that fraction, with a smoothing factor $\alpha$
setting how much of each new sample entered the estimate. The ablation
serves two purposes. It establishes that the bound reported in X2 is a
property of the design rather than of a filter setting, and it is the
measurement on which the decision to remove the filter rests.

The sweep covers $\alpha \in \{0.0, 0.1, \dots, 1.0\}$ at
$R \in \{1, 5, 10, 15\}$ across three seeds, 132 runs. The endpoints are
included deliberately: $\alpha = 0$ freezes the estimate and $\alpha = 1$
removes the smoothing entirely. The $\alpha = 1$ rows therefore describe
the oracle as it is evaluated in the rest of this chapter, and the
remaining rows describe what the filter was doing before it was taken
out.

Two caveats attach to the numbers below. They were collected on three
seeds rather than five, and against the version of the code that still
carried the filter, so they are not directly comparable digit-for-digit
with the five-seed tables above --- the $\alpha = 1$ column is the
like-for-like anchor, and it agrees with the X1 and X2 rows to within the
seed dispersion. The confirmation that matters does not rest on this sweep
alone: X1 and X2 were re-run after the removal, and the change they show
is reported in Section~\ref{sec:alpha-confirmed}.

\begin{table}[htbp]
  \centering
  \caption{Ablation: dirty-fraction and oracle-behaviour statistics by $\alpha$
    and $R$. The $\alpha = 0$ row is the degenerate case, included as a
    convergence check against pure BASE. Intermediate values of $\alpha$ are
    omitted here; the full grid is in the appendix.}
  \label{tab:alpha-grid}
  \small
  \begin{tabular}{r r r r r r r r}
    \toprule
    $R$ & $\alpha$ & $p_{50}$ (\%) & $p_{95}$ (\%) & Switches & $\phi_{\mathrm{time}}$ & $\phi_{\mathrm{comp}}$ & Thr.\ (tps) \\
    \midrule
    1   & 0.0      & 29.50         & 29.65         & 0.0      & 0.0000                 & 0.0000                 & 3.8654      \\
    1   & 0.1      & 5.97          & 11.11         & 57.7     & 0.9041                 & 0.1698                 & 0.0758      \\
    1   & 0.2      & 6.15          & 10.59         & 64.7     & 0.8654                 & 0.1341                 & 0.0737      \\
    1   & 0.5      & 6.49          & 10.22         & 74.0     & 0.8204                 & 0.1073                 & 0.0721      \\
    1   & 1.0      & 6.80          & 10.13         & 79.7     & 0.7866                 & 0.0924                 & 0.0720      \\
    \midrule
    5   & 0.0      & 27.20         & 27.53         & 0.0      & 0.0000                 & 0.0000                 & 3.3323      \\
    5   & 0.1      & 2.95          & 10.71         & 144.7    & 0.6656                 & 0.1280                 & 0.2787      \\
    5   & 0.2      & 4.43          & 10.53         & 177.3    & 0.5970                 & 0.1111                 & 0.3251      \\
    5   & 0.5      & 5.29          & 10.33         & 212.3    & 0.5491                 & 0.1011                 & 0.3410      \\
    5   & 1.0      & 5.67          & 10.25         & 232.7    & 0.5034                 & 0.0969                 & 0.3443      \\
    \midrule
    10  & 0.0      & 23.38         & 23.94         & 0.0      & 0.0000                 & 0.0000                 & 2.6651      \\
    10  & 0.1      & 0.65          & 10.11         & 157.7    & 0.5723                 & 0.0987                 & 0.4129      \\
    10  & 0.2      & 1.59          & 10.04         & 186.0    & 0.4391                 & 0.0841                 & 0.4922      \\
    10  & 0.5      & 2.39          & 9.99          & 206.7    & 0.3265                 & 0.0729                 & 0.5565      \\
    10  & 1.0      & 3.28          & 9.97          & 225.3    & 0.2920                 & 0.0702                 & 0.5863      \\
    \midrule
    15  & 0.0      & 17.74         & 18.46         & 0.0      & 0.0000                 & 0.0000                 & 1.9989      \\
    15  & 0.1      & 0.15          & 9.42          & 128.3    & 0.3947                 & 0.0847                 & 0.4997      \\
    15  & 0.2      & 0.29          & 9.39          & 146.3    & 0.2940                 & 0.0775                 & 0.5510      \\
    15  & 0.5      & 0.39          & 9.40          & 159.3    & 0.2032                 & 0.0726                 & 0.5835      \\
    15  & 1.0      & 0.43          & 9.34          & 160.7    & 0.1656                 & 0.0709                 & 0.5933      \\
    \bottomrule
  \end{tabular}
  \datasource{\texttt{experiments-old-ewma-0.2/alpha\_sweep/alpha\_sweep\_seed\{0,1,2\}.csv},
    132 runs, mean over 3 seeds.}
\end{table}

\subsection{The \texorpdfstring{$\alpha = 0$}{alpha = 0} anchor}\label{sec:alpha-anchor}

The degenerate setting is checked first, because it validates the
implementation. At $\alpha = 0$, across all four $R$ values and all three
seeds, the switch count, $\phi_{\mathrm{time}}$ and
\texttt{completed\_acid} are identically zero, not merely close. This is
the expected behaviour: with $\alpha = 0$ the smoothed estimate never
updates past its first sample, which is taken at an empty dirty set, so
the entry threshold is never crossed and the run is pure BASE. The
throughput of those rows agrees with the three-seed X1 BASE rows to five
significant figures at every $R$: 3.8654, 3.3323, 2.6651 and 1.9989. (The
five-seed BASE rows of Table~\ref{tab:x1-summary} differ from these in the
third decimal, which is the two additional seeds and not the oracle ---
BASE never consults it.) Together these confirm that $\alpha$ reached the
oracle rather than being silently ignored, and that the two sweeps agreed
where they overlapped.

\subsection{The bound does not depend on the smoothing}\label{sec:alpha-bound}

The main result of the sweep is negative in the useful sense: over the
whole engaged range of $\alpha$, from 0.1 to 1.0, the measured $p_{95}$
stays between 9.34\% and 11.11\% against the configured $c = 10\%$. The
entry side of the loop is therefore essentially decoupled from the
smoothing parameter, and the tracking result of
Section~\ref{sec:x2-critical-sweep} is not an artefact of any particular
choice of $\alpha$. This is also what makes removing the filter safe: the
guarantee the oracle is built to provide is the one quantity the filter
demonstrably does not touch.

The reason is the asymmetry described in
Section~\ref{sec:recovery-asymmetry}. Entry is evaluated in BASE, where
samples arrive every one to two simulated seconds and the dirty fraction
rises slowly, so even a heavily smoothed estimate crosses $c$ within a
fraction of a percentage point of the right moment. How much the estimate
is smoothed barely changes that.

\subsection{What the smoothing changed}\label{sec:alpha-effects}

Everything the smoothing factor affects sits on the exit side of the loop
and in the switching rate. Table~\ref{tab:alpha-removal} isolates the change
that removing it makes, by placing the earlier default of $\alpha = 0.2$
next to the $\alpha = 1.0$ rows that describe the oracle without a filter.

\begin{table}[htbp]
  \centering
  \caption{Ablation: the effect of removing the smoothing. $\alpha = 0.2$ is the
    earlier default, $\alpha = 1.0$ the unfiltered oracle.
    $\bar{d}_{\mathrm{exit}}$ is the mean live dirty fraction at the instant
    the oracle returns to BASE --- the effective recovery threshold of
    Section~\ref{sec:effective-recovery} --- against a configured
    $r = 5\%$.}
  \label{tab:alpha-removal}
  \small
  \begin{tabular}{r r r r r r}
    \toprule
    $R$ & $\alpha$ & $\bar{d}_{\mathrm{exit}}$ (\%) & $\phi_{\mathrm{time}}$ & Thr.\ (tps) & Switches \\
    \midrule
    1   & 0.2      & 2.01                           & 0.8654                 & 0.0737      & 64.7     \\
    1   & 1.0      & 3.57                           & 0.7866                 & 0.0720      & 79.7     \\
    \midrule
    5   & 0.2      & 0.79                           & 0.5970                 & 0.3251      & 177.3    \\
    5   & 1.0      & 3.46                           & 0.5034                 & 0.3443      & 232.7    \\
    \midrule
    10  & 0.2      & 0.21                           & 0.4391                 & 0.4922      & 186.0    \\
    10  & 1.0      & 2.84                           & 0.2920                 & 0.5863      & 225.3    \\
    \midrule
    15  & 0.2      & 0.03                           & 0.2940                 & 0.5510      & 146.3    \\
    15  & 1.0      & 1.98                           & 0.1656                 & 0.5933      & 160.7    \\
    \bottomrule
  \end{tabular}
  \datasource{\texttt{experiments-old-ewma-0.2/alpha\_sweep/alpha\_sweep\_seed\{0,1,2\}.csv},
    columns \texttt{exit\_dirty\_live\_mean}, \texttt{acid\_fraction\_time},
    \texttt{throughput\_tps} and \texttt{mode\_switch\_count}. Mean over 3
    seeds.}
\end{table}

\paragraph{The recovery threshold.} The filter is what made $r$
meaningless. Under smoothing the oracle leaves ACID with a true dirty
fraction of 2.01\% at $R = 1$ falling to 0.03\% at $R = 15$, against a
configured 5\%: the faster reconciliation runs, the further past the
configured level the system drains before the estimate catches up.
Removing the filter moves that exit point to 3.57\% and 1.98\%
respectively, and compresses its spread from $R = 1$ to $R = 15$ from a
factor of 67 to a factor of 1.8.

The unfiltered oracle still does not realise $r$ exactly, and the residual
gap is the point of Section~\ref{sec:effective-recovery}: with the filter
gone the oracle sees the true dirty fraction, but it only sees it when a
transaction arrives at admission, and between two arrivals a fast
reconciler pool can carry the backlog well past $r$. Removing the
smoothing eliminates one of the two lags; the sampling interval is the
other, and it is the one that remains.

\paragraph{ACID time and throughput.} Exiting earlier means spending less
of the run in ACID. $\phi_{\mathrm{time}}$ falls at every $R$, by 9\% in
relative terms at $R = 1$ and by 44\% at $R = 15$, and throughput moves
with it: from 0.325 to 0.344\,tps at $R = 5$, 0.492 to 0.586 at
$R = 10$ and 0.551 to 0.593 at $R = 15$. The one exception is $R = 1$,
where throughput is essentially flat and slightly lower, because at that
provisioning the oracle is in ACID for the great majority of the run under
either setting and the small change in $\phi_{\mathrm{time}}$ has little
left to recover. Since the unfiltered oracle also switches more often,
these gains rule out mode transitions as a significant direct cost: what
drives throughput is time spent in ACID, not the number of times the mode
changes.

\paragraph{Switching.} This is the one quantity that moves against the
removal. Switch counts rise from 64.7 to 79.7 at $R = 1$, 177.3 to 232.7
at $R = 5$, 186.0 to 225.3 at $R = 10$ and 146.3 to 160.7 at $R = 15$ ---
between 10\% and 31\% more transitions. Reducing this is the only thing
the filter measurably bought. The increase is bounded, though, and the
comparison against $\alpha = 0.1$ shows why: the hysteresis band, not the
filter, is doing most of the anti-flap work, and the filter was
contributing a further reduction on top of it rather than providing the
mechanism.

\subsection{Why the smoothing was removed}\label{sec:alpha-interpretation}

The sweep separates the two halves of the oracle. The entry side, which
carries the guarantee, is insensitive to the smoothing factor over its
entire range. The exit side, which determines how much throughput is given
up and whether $r$ names a real operating point, is sensitive to it
throughout.

That split is what decided the design. Filtering is worth its cost when a
signal carries measurement noise, and the dirty fraction carries none: it
is an exact count of a set the simulator maintains, not an estimate of an
unobservable quantity. What the filter attenuated was therefore not noise
but real movement, and it did so on top of a second lag --- the
admission-driven sampling interval --- that cannot be removed by tuning.
Of the two, the filter was the larger contributor on the exit side, and
the two together drove the effective recovery threshold two orders of
magnitude below its configured value at $R = 15$.

Set against that, the filter bought a 10--31\% reduction in the switching
rate, on a mechanism the hysteresis band already provides. Removing it
leaves the guarantee untouched, returns between a tenth and a half of the
ACID time budget depending on provisioning, and makes the recovery
threshold mean something much closer to what it says. The oracle described
in Section~\ref{sec:approach-adaptive-oracle} therefore acts on the
instantaneous dirty fraction, and the smoothing factor is no longer a
parameter of the system.

\subsection{Confirmed by re-running the sweeps}\label{sec:alpha-confirmed}

The sweep above predicted the removal would be safe. X1 and X2 were then
re-run against the filter-free implementation on five seeds, which turns
that prediction into a measurement. Table~\ref{tab:alpha-confirmed} places
the two sets of \texttt{HYBRID\_ADAPTIVE} rows side by side at the
baseline $c = 0.10$.

\begin{table}[htbp]
  \centering
  \caption{Effect of removing the smoothing, measured by re-running X1.
    ``before'' is the filtered oracle at the earlier default
    $\alpha = 0.2$, three seeds; ``after'' is the filter-free oracle, five
    seeds. The guarantee ($p_{95}$ against $c = 10\%$) is unchanged; the
    cost of providing it falls.}
  \label{tab:alpha-confirmed}
  \small
  \begin{tabular}{r r r r r r r r r}
    \toprule
        & \multicolumn{3}{c}{Thr.\ (tps)} & \multicolumn{3}{c}{$\phi_{\mathrm{time}}$} & \multicolumn{2}{c}{$p_{95}$ (\%)}                                                \\
    \cmidrule(lr){2-4}\cmidrule(lr){5-7}\cmidrule(lr){8-9}
    $R$ & before                          & after                                      & change                            & before & after  & change    & before & after \\
    \midrule
    1   & 0.0737                          & 0.0726                                     & $-1.5$\%                          & 0.8654 & 0.7776 & $-10.1$\% & 10.59  & 10.12 \\
    5   & 0.3251                          & 0.3453                                     & $+6.2$\%                          & 0.5970 & 0.5068 & $-15.1$\% & 10.53  & 10.22 \\
    10  & 0.4922                          & 0.5883                                     & $+19.5$\%                         & 0.4391 & 0.2886 & $-34.3$\% & 10.04  & 9.98  \\
    15  & 0.5510                          & 0.5941                                     & $+7.8$\%                          & 0.2940 & 0.1667 & $-43.3$\% & 9.39   & 9.36  \\
    \bottomrule
  \end{tabular}
  \datasource{``before'' from
    \texttt{experiments-old-ewma-0.2/reconciler\_sweep/reconciler\_sweep\_seed\{0,1,2\}.csv}
    at $\alpha = 0.2$; ``after'' from Table~\ref{tab:x1-summary}.}
\end{table}

The measured outcome is slightly better than the sweep predicted. The
$p_{95}$ column moves by at most 0.47\,pp and stays inside the X2 tracking
bound at every $R$, so the guarantee is untouched --- and across the full
X2 grid that bound in fact tightened, from 1.3\,pp to 0.74\,pp. The ACID
time share falls at every $R$, by 10\% at $R = 1$ and 43\% at $R = 15$,
and throughput rises at every $R$ except the first, by up to 19.5\%. The
switching rate rises correspondingly, from 65 to 80 switches at $R = 1$
and from 146 to 162 at $R = 15$, which is the cost the hysteresis band was
always paying for and which the filter was quietly reducing on top of it.

The largest single effect is the one reported in
Section~\ref{sec:effective-recovery}: the effective recovery threshold. Under
the filter the mean live dirty fraction at exit fell from 2.01\% at
$R = 1$ to 0.03\% at $R = 15$; without it, it runs 3.64\% to 2.05\%
against a configured $r$ of 5\%. The configured recovery threshold went
from describing nothing at high $R$ to describing an upper bound that is
approached within 3\,pp. Removing the filter therefore did not trade the
guarantee for throughput --- it returned throughput, ACID budget and the
meaning of $r$ at once, which is what should be expected when a filter is
applied to a signal that carries no noise.

\section{The effective recovery threshold}\label{sec:effective-recovery}

Every hybrid run in this chapter carries a second oracle threshold that no
sweep varies: the oracle enters ACID when the dirty fraction reaches the
critical threshold $c$, and returns to BASE once it has fallen back to the
recovery threshold $r$ --- $r = 0.05$ at the baseline, $r = 0.5c$
throughout X2. The effect reported here is present in all five sweeps,
which is why it is placed after them rather than inside one of them. In
one sentence: $r$ is an upper bound on where an ACID excursion ends, not
the point at which it ends, and runs leave ACID cleaner than they were
asked to --- by more, the faster reconciliation is provisioned.

\begin{table}[htbp]
  \centering
  \caption{Where the oracle actually switches, how long an ACID transaction
    waits, and how often the oracle looks. Measured under
    \texttt{HYBRID\_ADAPTIVE} at $c = 0.10$, $r = 0.05$. Dirty fractions
    are percentages of the 5000 resources; the \emph{miss} columns are the
    distance from the threshold that fired, in percentage points, positive
    meaning dirtier than configured.}
  \label{tab:effective-recovery}
  \label{tab:sampling-interval}
  \small
  \setlength{\tabcolsep}{4pt}
  \begin{tabular}{r r r r r r r r r}
    \toprule
        & \multicolumn{2}{c}{Entry (BASE $\to$ ACID)} & \multicolumn{3}{c}{Exit (ACID $\to$ BASE)} & Mean ACID & \multicolumn{2}{c}{Interval between looks (s)}                                           \\
    \cmidrule(lr){2-3}\cmidrule(lr){4-6}\cmidrule(lr){8-9}
    $R$ & mean $d$                                    & miss vs $c$                                & mean $d$  & miss vs $r$                                    & $p_{95}$ & wait (s) & in ACID & in BASE \\
    \midrule
    1   & 10.08                                       & $+0.08$                                    & 3.64      & $-1.36$                                        & 4.88     & 392.5    & 107.2   & 3.37    \\
    5   & 10.11                                       & $+0.11$                                    & 3.42      & $-1.58$                                        & 4.92     & 70.2     & 15.0    & 1.58    \\
    10  & 10.16                                       & $+0.16$                                    & 2.85      & $-2.15$                                        & 4.86     & 32.5     & 7.0     & 1.30    \\
    15  & 10.12                                       & $+0.12$                                    & 2.05      & $-2.95$                                        & 4.74     & 22.5     & 3.9     & 1.51    \\
    \bottomrule
  \end{tabular}
  \datasource{\texttt{critical\_sweep\_seed\{0..4\}.csv}, $c = 0.10$ rows;
    columns \texttt{entry\_dirty\_live\_mean},
    \texttt{exit\_dirty\_live\_\{mean,p95\}} and
    \texttt{acid\_dirty\_block\_time\_mean\_ms}, mean over 5 seeds. The
    transitions are counted from \texttt{entry\_dirty\_live\_count}: 40,
    117, 112 and 81 entries respectively. \emph{Mean ACID wait} is the
    average length of one wait on a repair, per wait rather than per
    transaction, as in Table~\ref{tab:blocking}. The intervals are
    simulated time in each mode divided by the root transactions tagged in
    that mode; \texttt{started\_\{acid,base\}} counts those tags.}
\end{table}

\paragraph{What a look is.} The last two columns need saying plainly,
because the oracle has no timer and the word \emph{interval} would
otherwise suggest one. There is exactly one place where the rule is
evaluated: a root transaction reaches admission carrying no mode tag, and
the routine that gives it one reads the live dirty count and compares it
against $c$ and $r$ (Section~\ref{sec:approach-sampling}). One tagged root
transaction is one look, and nothing else --- not the passing of time, not
a bookkeeper transaction, not a reconciliation finishing --- causes
another. Between two looks the oracle simply does not know what the dirty
set is doing. The two columns divide the simulated time spent in each mode
by the number of transactions tagged in it, so they report how long the
rule went unevaluated on average: every 1.3 to 3.4 seconds in BASE, and
every 3.9 to 107 seconds in ACID.

This is the oracle's own sensor and not the thesis's measuring instrument.
The dirty fractions reported everywhere else in this chapter come from the
metrics sampler, which fires once per simulated second in both modes
(Section~\ref{sec:exp-metrics}); it is the divergence between the two that
this section is about.

\paragraph{Entry is prompt, exit is late.} At entry the threshold means
what it says: the oracle fires 0.08 to 0.16\,pp above the configured 10\%,
with no trend in $R$. Exit misses, and always on the clean side. The mean
dirty fraction at the return to BASE is 3.64\% at $R = 1$, which is
1.36\,pp below the configured $r$ of 5\%, and 2.05\% at $R = 15$, which is
2.95\,pp below it; the miss grows steadily with reconciliation capacity.
Figure~\ref{fig:recovery-gap} is the two edges drawn against the
thresholds they were given. The $p_{95}$ column caps the miss: the
dirtiest exits sit at 4.74--4.92\%, just under $r$ at every $R$.

\begin{figure}[htbp]
  \centering
  \includegraphics[width=\textwidth]{recovery_gap.pdf}
  \caption{What the thresholds were set to, and where the oracle actually
    switched, at $c = 0.10$ and $r = 0.05$. Entry lands on $c$ at every
    reconciler count; exit falls short of $r$ by a distance that grows with
    reconciliation capacity. The short bars are the $p_{95}$ exit, the
    dirtiest observed. Means over five seeds, the data of
    Table~\ref{tab:effective-recovery}.}
  \label{fig:recovery-gap}
\end{figure}

\paragraph{Why the exit runs past $r$.}\label{sec:recovery-mechanism}
Looks are scarce in ACID because ACID is the mode in which transactions
stop arriving. One that needs a resource the reconcilers have not cleared
is parked until they do (Section~\ref{sec:approach-arbiter}), it holds its
in-flight slot while parked, and under \texttt{EQUILIBRIUM} nothing new is
admitted until something completes. The wait column is the size of that
stall: one wait on a repair averages 22.5\,s at $R = 15$ and 392\,s at
$R = 1$, against mean excursions of 103\,s and 976\,s. Hence the 3.9 to
107\,s between looks in ACID against 1.3 to 3.4\,s in BASE --- and a
larger budget lengthens the wait further, because it grows with the
dirtiness the transaction was admitted into (at $R = 10$, 7.9\,s at
$c = 0.02$ against 59.0\,s at $c = 0.20$, Table~\ref{tab:blocking}).

The interval is far shorter than the wait, and why it is says what the
interval column hides. The thirty in-flight transactions are not all
parked at the same moment; the ones that are not still finish and free
their slots, so admissions keep trickling in and the average comes out at
a few seconds even where a single wait runs to twenty of them. What the
average smooths over is the stretches in which enough of the population is
parked at once that nothing is admitted at all, and it is those, not the
mean, that carry an exit past $r$. One excursion at $R = 15$, whose mean
interval is the 3.9\,s of the table, contains a 33-second run in which the
oracle is never consulted; the dirty set crosses $r$ inside it and is at
1.66\% by the time the next transaction arrives to notice.

Repair meanwhile neither pauses nor draws on
the foreground pool: under pure BASE the reconcilers complete exactly
$25{,}000R$ reconciliations for every $R \le 20$, a drain of $0.5R$
resources per second, against a standing backlog of 394 to 1473 resources
that Little's law matches to the measured recovery latency within 3\%.
Only at $R = 25$ does the pool outrun the workload, its backlog collapsing
to 2.1 resources, which is what makes $R = 25$ degenerate in X1
(Section~\ref{sec:x1-engagement}). The dirty set therefore keeps draining
at a rate set by $R$ while the foreground is stalled and the oracle is not
looking, and the excursion ends wherever the drain has got to when some
transaction finally reaches admission. That gives the direction of the
effect but not its size: the mean interval times the drain rate accounts
for the undershoot at $R = 1$ and is too small at $R = 15$, and the
distribution of the looks within an excursion, which is what would settle
it, is not measured here.


\paragraph{The slack protects the budget rather than weakening
  it.}\label{sec:recovery-asymmetry} The coarse sampling does not reach
the two edges equally. Entry is decided in BASE, where a look arrives
every one to three seconds and dirtiness accumulates slowly under
stationary load; exit is decided in ACID, where looks are rare and the
quantity moves fastest. Exits are therefore late and never early, and a
run ends up cleaner than its configuration implies rather than dirtier, so
the tail bound of X2 is unaffected by the slack in $r$
(Section~\ref{sec:x2-tracking}). What it costs is ACID time rather than
consistency: excursions run 2.0 to 3.1 times longer than the drain from
$cN = 500$ to $rN = 250$ resources requires --- 976\,s against 500\,s at
$R = 1$, 103\,s against 33\,s at $R = 15$. Sampling on a fixed interval of
simulated time rather than once per admission would remove most of that,
and is left to future work.


\section{Threats to validity}\label{sec:exp-threats}

\begin{itemize}
  \item \textbf{Stationary load only (\texttt{EQUILIBRIUM}).} Adaptive
        control would be expected to show its clearest advantage when the
        operating point moves, and constant load cannot test that. The claim is
        scoped accordingly in Section~\ref{sec:approach-why-hybrid}, and the X3
        claim about the knob (O9) is phrased in terms of transfer across
        \emph{configurations} rather than adaptation within a run. The
        stationarity of the load is stated at the head of
        Section~\ref{sec:x3-random-baseline} as well, because the throughput
        comparison there is run in the regime least favourable to feedback and
        should be read in that light.

  \item \textbf{Five seeds per configuration.} Mitigated but not removed by
        how little the results move between seeds. The random arm of X3 is the
        high-variance one --- 5.52\% median seed CV on $p_{95}$ against 0.25\%
        for the adaptive arm --- so the O9 conclusion rests on the arm
        with the larger error bars, in the direction that favours the
        conclusion. The throughput ratios of Table~\ref{tab:x3-matched} carry
        seed standard deviations of 0.004 to 0.051, which is small against the
        0.82 to 1.39 range of the ratios themselves but is computed from five
        runs and should be read as an indication rather than an interval.
        The ablation of Section~\ref{sec:alpha-ablation} is the exception at three
        seeds, and its conclusion is independently confirmed at five in
        Section~\ref{sec:alpha-confirmed}.

  \item \textbf{The blocking counter is not a transaction count.}
        \texttt{acid\_txn\_blocked\_on\_dirty\_count} counts distinct root
        transaction identifiers that blocked at least once, and the set it
        counts is not a subset of the admitted one --- it includes
        transactions still parked at the end of the run and transactions that
        aborted while parked. It is therefore reported in this chapter only as
        a raw activity total (Table~\ref{tab:x0-arbiter}) and never as a rate
        against admissions, and no claim rests on it. Measuring how often an
        ACID transaction blocks, rather than how long it waits, would need a
        per-transaction instrument.

  \item \textbf{\texttt{reconciler\_busy\_fraction} is misnamed.} It is not a
        utilisation but the fraction of samples in which \emph{every}
        reconciler was busy at the same instant, taken once per simulated
        second against a service time of two seconds. The phase relationship
        between the two makes it unreliable: it reads 0.5 at $R = 1$ and 0.0
        at $R \ge 10$ under pure BASE, while the same runs sustain exactly
        $0.5R$ completed reconciliations per second against a standing queue
        of over a thousand resources --- a pool that is in fact fully
        saturated. Saturation is established in this chapter from the
        completion rate and the queue depth instead
        (Section~\ref{sec:recovery-mechanism}); no result uses this column,
        and it should not be read from the exported CSVs.

  \item \textbf{$R = 0$ is a different mechanism, not zero capacity.} At
        $R = 0$ the CPR work is performed inline by the admitting transaction,
        so the run has no reconciliation queue and 1.16\,M mark events against
        26\,k at $R = 1$. The $R = 0$ rows are a reference line and are not
        points on the same capacity axis as $R \ge 1$.

  \item \textbf{Simulated rather than real execution, and a single workload
          shape.} The absolute numbers are properties of the simulator; the
        comparisons between configurations are the transferable part.

  \item \textbf{The matched comparison of X3 reads between measured
          points.} The random arm is sampled at five values of $p$, and
        Table~\ref{tab:x3-matched} reads its throughput at intermediate
        values of $p_{95}$ rather than at measured ones. Only four of the
        thirteen comparisons land within 15\% of a measured random run; the
        rest lean on the estimate, and recomputing them against the nearest
        measured run instead moves several by more than a tenth and reverses
        the sign of one. No conclusion is drawn from those rows. The interval
        $p \in (0, 0.05)$ is unsampled and contains the steepest part of the
        trade-off; a finer grid there, together with settings between
        $p = 0.10$ and $p = 0.25$ to bracket the tight-budget adaptive cells
        and settings chosen to reach the dirty end at high $R$, would turn
        most of these readings into measurements and close the six cells
        where the two arms do not overlap at all.

  \item \textbf{The workload sweeps vary one axis at a time.} X4 and X5
        establish that the budget holds across semantic overlap and across
        resources per transaction, but each holds the other at its
        baseline. The corner where
        both are extreme is unsampled, and X4's saturation result at
        $\sigma = 1.00$, $R = 1$ suggests it is the corner most likely to
        behave differently.

  \item \textbf{$L = 400$ leaves the regime under test.} At that setting the
        30 transactions held in flight demand more resources between them than
        the system has locks, so lock conflict rather than the oracle limits
        admission and the abort rate reaches 0.65. The block is a stress anchor, not a scaling point,
        and is excluded from the invariance claim of
        Section~\ref{sec:x5-invariance}.

  \item \textbf{The ablation cannot be re-run against current code.} The
        \texttt{alpha\_sweep} results behind Section~\ref{sec:alpha-ablation}
        are archived under \texttt{experiments-old-ewma-0.2/} and the tables
        can be regenerated from them, but the parameter they sweep no longer
        exists in the simulator, so the runs themselves cannot be reproduced
        without reverting the removal. The confirmation in
        Section~\ref{sec:alpha-confirmed} is reproducible against current
        code; the ablation grid is reproducible only from the archive.
\end{itemize}

% NOTE: moved here from Section 5.1 pending a decision on placement.
\section{The earlier filtered runs}\label{sec:exp-earlier-runs}

Two generations of results exist for this system, and the artefact archive
contains both, so it is worth saying at the outset which is which and why
the older set was superseded.

Every number in this chapter, unless the text says otherwise, comes from
the oracle described in Section~\ref{sec:approach-adaptive-oracle}, which
compares its thresholds against the dirty fraction read at the instant of
the decision. An earlier version instead compared them against an
exponentially weighted moving average of that fraction,
\begin{equation}
  \hat{d}_k = \alpha d_k + (1-\alpha)\hat{d}_{k-1},
  \label{eq:ewma}
\end{equation}
where $\alpha \in [0,1]$ is the smoothing factor and $k$ indexes
successive samples; larger $\alpha$ tracks the signal more closely and
filters less. The default was $\alpha = 0.2$. One property of the filter
matters for what follows: its time constant is expressed in \emph{samples},
whereas the quantity it estimates evolves in \emph{seconds}, and because
the oracle samples once per admission rather than once per unit of
simulated time (Section~\ref{sec:approach-sampling}), the rate relating
the two depends on load.

X1, X2 and X3 were run under both, on three seeds under the filtered
oracle and five under the current one; those older files are archived
under \texttt{experiments-old-ewma-0.2/}.

The filter was introduced as an anti-flap measure, and the reason it was
removed is that it did not work as one. Table~\ref{tab:ewma-note}
decomposes its effect. It is true that the filter lowered the raw switch
count, by 9 to 24\%, and read alone that column supports the
anti-flap story. The decomposition does not. Switches fall for two quite
different reasons --- because the system enters ACID less often, or
because each visit lasts longer --- and only the first is stabilisation.

\begin{table}[htbp]
  \centering
  \caption{What the smoothing filter actually did. \emph{Entries} counts
    BASE $\to$ ACID transitions and \emph{BASE dwell} is mean simulated
    time in BASE per entry, that is, how long the system stayed out of
    ACID between excursions. $\alpha = 0.2$ is the earlier default,
    $\alpha = 1.0$ the same code path with the filter disabled, so the two
    columns differ in nothing else.}
  \label{tab:ewma-note}
  \small
  \begin{tabular}{r r r r r r r r r}
    \toprule
        & \multicolumn{2}{c}{Entries} & \multicolumn{3}{c}{BASE dwell per entry (s)} & \multicolumn{3}{c}{Excursion length (s)}                                                                        \\
    \cmidrule(lr){2-3}\cmidrule(lr){4-6}\cmidrule(lr){7-9}
    $R$ & $\alpha{=}0.2$              & $\alpha{=}1.0$                               & $\alpha{=}0.2$                           & $\alpha{=}1.0$ & change  & $\alpha{=}0.2$ & $\alpha{=}1.0$ & change  \\
    \midrule
    1   & 33                          & 40                                           & 206                                      & 264            & $-22$\% & 1330           & 987            & $+35$\% \\
    5   & 89                          & 117                                          & 227                                      & 213            & $+7$\%  & 336            & 217            & $+55$\% \\
    10  & 93                          & 113                                          & 302                                      & 314            & $-4$\%  & 236            & 130            & $+82$\% \\
    15  & 73                          & 81                                           & 481                                      & 517            & $-7$\%  & 201            & 103            & $+95$\% \\
    \bottomrule
  \end{tabular}
  \datasource{\texttt{experiments-old-ewma-0.2/alpha\_sweep/alpha\_sweep\_seed\{0,1,2\}.csv},
    $c = 0.10$, mean over 3 seeds. ``Change'' is the filtered value
    relative to the unfiltered one. BASE dwell is
    $(1-\phi_{\mathrm{time}})\times$ duration divided by
    \texttt{entry\_dirty\_live\_count}.}
\end{table}

The BASE dwell columns are the ones that answer the question. At
$R = 5, 10, 15$ the filter changes how long the system stays out of ACID
by between $-7$\% and $+7$\% --- that is, not at all, to within the seed
dispersion. The rate at which the system decides it needs ACID is set by
the entry threshold and the workload, and the filter does not touch it.
What the filter changes is the other term: each excursion runs 35 to 95\%
longer, because the smoothed estimate is still above $r$ well after the
dirty set has drained below it. The fall in switch count is arithmetic
--- fewer, longer visits in a fixed run length --- and not stability.

The cost of buying that arithmetic is paid in the mode the system is stuck
in: ACID time share rises by 10\% at $R = 1$ and 77\% at $R = 15$
relative to the unfiltered oracle. The filter did not damp an oscillation;
it added lag on the exit edge, and the reduced switch count was a symptom
of the lag rather than a benefit.
Section~\ref{sec:alpha-ablation} reports the full sweep of $\alpha$ behind
this, and Section~\ref{sec:alpha-confirmed} the effect of the removal on
X1 and X2.

